\documentclass[11pt]{article}
\usepackage[T1]{fontenc}
\usepackage[margin=1in]{geometry}
\usepackage{tikz}
\usepackage{pgfplots}
\pgfplotsset{compat=1.18}
\usepackage[utf8]{inputenc}
\usepackage{booktabs}
\usepackage{graphicx}
\usepackage{amsmath}
\usepackage{amssymb}
\usepackage[version=4]{mhchem}
\usepackage{siunitx}
\usepackage{longtable,tabularx}
\setlength\LTleft{0pt} 
\usepackage{xcolor}
\usepackage{svg}
\usepackage{float}
\usepackage{placeins}
\usepackage{ifthen}
\usepackage{natbib}
\usepackage{url}
\usepackage{hyperref}
\svgsetup{inkscapepath=fig/}
\svgpath{{fig/}}  % directory containing figures
\hypersetup{hidelinks}
\setcounter{MaxMatrixCols}{20}

\newboolean{commentsOn}
\setboolean{commentsOn}{true} % or false

\begin{document}

\begin{center}
{\LARGE\bfseries Aircraft System Identification Benchmark Repository Reference}\\[0.75em]
{\large Classical, Grey-Box, Surrogate, and Learning-Based Methods for 3-DOF and 6-DOF Aircraft Data}\\[0.75em]
{\normalsize Generated reference document for the \texttt{system\_identification} repository}\\[0.5em]
{\normalsize \today}
\end{center}
\vspace{1em}

\begin{abstract}
This document is the technical reference for the \texttt{system\_identification}
repository. It defines the benchmark motivation, aircraft models, generated
datasets, measurement assumptions, hidden-control cases, implemented
identification methods, code-to-equation traceability, validation metrics, and
current result summaries. The repository focuses on aircraft system
identification for small fixed-wing vehicles when motion capture may be the
only trusted measurement source. The current benchmark suite includes a
nonlinear 3-DOF longitudinal model with nominal-plus-unmodeled aerodynamic
terms, direct-state and motion-capture-derived observation modes, hidden
SAFE-like command reshaping, multiple maneuver families, and a developing
6-DOF model/data interface. The purpose of this reference is reproducibility:
it should be possible to trace each equation, dataset field, result table, and
figure back to repository code and regenerated artifacts.
\end{abstract}

\tableofcontents
\clearpage

\section*{Nomenclature}

\begin{longtable}{@{}l @{\quad=\quad} l@{}}
$N$ & number of discrete samples \\
$n_x, n_u, n_y$ & dimensions of state, input, and output vectors \\
$i$ & discrete-time index ($i = 1, \ldots, N$) \\[3pt]
$\mathbf{x}(t) \in \mathbb{R}^{n_x}$ & continuous-time state vector \\
$\mathbf{u}(t) \in \mathbb{R}^{n_u}$ & continuous-time control input \\
$\mathbf{y}(t) \in \mathbb{R}^{n_y}$ & continuous-time measured output \\
% $\boldsymbol{\theta}$ & vector of unknown model parameters \\
$p_x,p_z$ & inertial longitudinal and vertical position measured by motion capture \\
$V$ & true airspeed (wind axis) \\
$\alpha$ & angle of attack\\
$\gamma$ & flight–path angle \\
$Q$ & body pitch rate about $y$–axis \\
$\theta$ & pitch attitude \\
${\theta_p}$ & vector of constant model parameters \\
$T$ & thrust along $b_x$ \\
$\delta_e$ & elevator deflection \\
$L, D, M$ & lift, drag, and pitching moment \\
$m, J_y$ & mass and pitch moment of inertia \\
$\rho, S$ & air density and reference area \\
$g$ & gravitational acceleration \\
\end{longtable}


\section{Repository Scope and Orientation}

This reference documents the benchmark implemented in the \texttt{system\_identification} repository. It is not a general literature survey and it is not formatted as a journal submission. The purpose is to make the repository auditable: dataset definitions, model equations, method assumptions, generated figures, and result tables should all trace back to code that can be rerun.

The current repository has three main products:
\begin{enumerate}
    \item a reproducible 3-DOF longitudinal aircraft benchmark with direct-state and motion-capture-style observation channels;
    \item a nonlinear 6-DOF aircraft dataset with stall/residual aerodynamics for future coupled-dynamics benchmarks;
    \item a static benchmark website data export and method-plugin interface for future contributed methods.
\end{enumerate}

The primary motivation is small fixed-wing system identification in indoor test environments such as the PURT pylon-racing and fixed-wing UAS testbed setting described by \citet{sribunma2026indoor}. In that setting, the trusted measurement source may be external motion capture rather than onboard air-data or inertial sensors. The 3-DOF benchmark therefore reports both idealized direct-state results and motion-capture-derived results, with additional mocap-output estimators for methods that fit position and attitude directly instead of differentiating them into aerodynamic states.

\section{Repository Method Inventory}

This section lists only the method families currently represented in repository code or generated benchmark tables. The implementations are centered in \texttt{methods/comparison\_suite.py}, with the emerging public plugin interface under \texttt{methods/plugins/}. Each method is evaluated by training on its assigned dataset and then performing validation rollouts using the validation initial condition and pilot-command input history. Validation measurements are used for scoring, not for feedback during the rollout, unless a method is explicitly categorized as a measurement-assimilating estimator.

\subsection{Baselines and Fixed-Structure Aircraft Identification}

\paragraph{Nominal model.}
The nominal row evaluates the coded 3-DOF aircraft model without the hidden nonlinear aerodynamic residuals. It is a no-fit baseline and appears in the result figures as the reference performance level.

\paragraph{EquationError-LS.}
The equation-error least-squares row fits aerodynamic coefficients from reconstructed state derivatives. It is fast and interpretable, but it is sensitive to measurement noise because the dependent variables require numerical differentiation.

\paragraph{OEM-SS.}
The output-error single-shooting row estimates fixed aerodynamic parameters by optimizing trajectory error through the longitudinal model. The implementation uses CasADi/IPOPT and is the main deterministic fixed-structure aircraft identification baseline for direct-state data.

\paragraph{OEM-MocapOutput.}
The mocap-output OEM row treats $V,\alpha,\gamma,Q$ as latent states and fits measured position and pitch attitude directly. This is the most relevant deterministic baseline for the motion-capture-only setting because it avoids differentiating position and attitude to construct aerodynamic states.

\paragraph{Variational-Mocap.}
The variational/multiple-shooting mocap row adds process-residual flexibility while estimating latent states from mocap outputs. It is included as a practical bridge between fixed-structure output error and process-noise or smoothing formulations.

\paragraph{EKF-ParamID and Fisher-UQ.}
The EKF parameter-identification row estimates parameters from training data and then validates the learned model in open-loop rollout. The Fisher-information row provides local uncertainty and parameter-correlation diagnostics for the fitted fixed-structure model.

\subsection{Linear, Local, Lifted, and Frequency-Domain Predictors}

\paragraph{Linear-SS.}
The linear state-space row fits a discrete affine state-space predictor. It is a strong local predictive baseline, especially near trim or when mocap-derived state reconstruction dominates the error budget, but it should not be interpreted as recovering nonlinear aerodynamic structure.

\paragraph{Model-Stitching.}
The model-stitching row fits local linear models on trim-grid data and blends them using scheduling weights. It tests whether multiple local models improve over one global linear model across wider initial-condition and maneuver variation.

\paragraph{Subspace-Hankel.}
The subspace row fits a compact linear predictor using Hankel/subspace-style structure. It represents the repository's current state-space realization baseline.

\paragraph{Koopman-EDMD.}
The Koopman-EDMD row fits a lifted linear predictor using a fixed nonlinear feature library. It is included as a predictive surrogate rather than as an aerodynamic-parameter estimator.

\paragraph{Frequency-Welch.}
The frequency-domain row uses Welch/ETFE-inspired spectral estimates, coherence screening, and a weighted continuous-time linear fit. It is not CIFER and should be interpreted only as the repository's compact spectral baseline.

\paragraph{Frequency-Stitching.}
The frequency-stitching row applies the same spectral fitting idea locally on trim-grid groups and blends the resulting local continuous-time models. It tests whether the poor global frequency-domain performance is primarily a single-LTI-model limitation.

\subsection{Sparse, Symbolic, Kernel, and Neural Closure Methods}

\paragraph{SINDy.}
The SINDy row uses a structured coefficient-library formulation so that known low-order aircraft terms can be protected while sparse nonlinear residual terms are selected from a candidate library.

\paragraph{Symbolic-Stepwise.}
The symbolic-stepwise row uses the same candidate-library idea with forward stepwise selection. It is a lightweight symbolic baseline rather than a full genetic-programming symbolic-regression implementation.

\paragraph{GP-CoeffClosure.}
The GP/RBF coefficient-closure row learns a smooth residual aerodynamic coefficient map with kernel features. It represents the repository's current kernel surrogate closure method.

\paragraph{NN-CoeffSurrogate.}
The neural coefficient surrogate row learns aerodynamic coefficient residuals from labels available in the synthetic benchmark. Because real flight tests do not provide those coefficient labels directly, this row is an upper-bound diagnostic for what a supervised closure could achieve.

\paragraph{PINN-CoeffClosure and PINN-HiddenElevator.}
The PINN-style rows use flight-dynamics-level residuals, not full CFD or Navier--Stokes constraints. The coefficient-closure row learns aerodynamic closures, while the hidden-elevator row tests whether a neural model can infer a hidden command correction in SAFE-like cases.

\paragraph{UDE-Residual and UDE-HiddenControl.}
The UDE rows retain the known longitudinal dynamics and learn residual terms. The residual row targets missing aerodynamics, while the hidden-control row targets unknown input reshaping from a proprietary SAFE-like controller.

\subsection{6-DOF Baseline Methods}

\paragraph{6DOF-Nominal.}
The 6-DOF nominal row propagates the repository's attached-flow quaternion rigid-body model with the supplied pilot-command history. It has no fitted parameters and is used as the no-fit baseline against the hidden nonlinear stall/residual aerodynamic truth model.

\paragraph{6DOF-LinearSS.}
The 6-DOF linear state-space row fits an affine one-step predictor,
\[
\mathbf{x}_{k+1}=A\mathbf{x}_k+B\mathbf{u}_{\mathrm{cmd},k}+\mathbf{c},
\]
using ridge-regularized least squares. During validation the model is rolled out open-loop from the validation initial condition and pilot-command sequence.

\paragraph{6DOF-RidgeResidual.}
The 6-DOF residual row first propagates the nominal RK4 model and then adds a ridge-fitted one-step residual correction,
\[
\mathbf{x}_{k+1}=\Phi_{\mathrm{nom}}(\mathbf{x}_k,\mathbf{u}_{\mathrm{cmd},k})
+W^T[\mathbf{x}_k^T,\mathbf{u}_{\mathrm{cmd},k}^T,1]^T.
\]
This tests the same grey-box idea used in UDE-style methods, but with a linear residual so that the 6-DOF benchmark remains fast and deterministic. The residual absorbs actuator lag and the hidden difference between the attached-flow nominal model and the stalled nonlinear truth model.

\paragraph{6DOF-OEM-MocapOutput.}
The mocap-output 6-DOF row fits an affine one-step predictor on measured position and quaternion attitude. It does not predict body velocity or angular rate, so its validation score is a mocap-output NRMSE rather than a full-state NRMSE.

\section{Dataset and Result Artifacts}

The 3-DOF benchmark datasets are generated under \texttt{methods/data/longitudinal\_3dof\_nonlinear\_*}. The generated splits contain direct state measurements, mocap position/attitude measurements, mocap-derived aerodynamic states, pilot commands, actuator-realized commands, hidden SAFE corrections, coefficient residual labels, and residual-dynamics labels. The default validation metric in the result tables is the mean range-normalized state RMSE, denoted $J_{\mathrm{val}}$ in this document.

The 6-DOF dataset interface is generated by
\begin{verbatim}
./results.py simulate-6dof
\end{verbatim}
and currently writes \texttt{methods/data/aircraft\_6dof\_mixed/}. The full local 6-DOF workflow is
\begin{verbatim}
./results.py all-6dof
\end{verbatim}
which generates train/validation data, runs the implemented 6-DOF baseline methods, exports website data, and refreshes LaTeX-ready artifacts. The 6-DOF dataset contains inertial position, body velocity, quaternion attitude, body angular rates, pilot commands, actuator-realized commands, noisy direct-state measurements, mocap-style position/quaternion measurements, true aerodynamic coefficients, attached-flow nominal coefficients, and coefficient residual labels.

\section{6-DOF Coupled-Dynamics Benchmark}

The 6-DOF benchmark is the repository's nonlinear coupled-dynamics extension. The state is
\[
\mathbf{x}=
\begin{bmatrix}
x_n & y_e & z_d & u & v & w & q_w & q_x & q_y & q_z & p & q & r
\end{bmatrix}^T,
\]
where $(x_n,y_e,z_d)$ is inertial position, $(u,v,w)$ is body velocity, $(q_w,q_x,q_y,q_z)$ is the unit quaternion mapping body coordinates to inertial coordinates, and $(p,q,r)$ are body angular rates. The input is
\[
\mathbf{u}_{\mathrm{cmd}}=
\begin{bmatrix}
\delta_T & \delta_e & \delta_a & \delta_r
\end{bmatrix}^T.
\]
The simulator applies first-order actuator lag to obtain $\mathbf{u}_{\mathrm{act}}$, but the identification methods are supplied $\mathbf{u}_{\mathrm{cmd}}$ to preserve the practical pilot-command setting.

The rigid-body equations use body-axis velocity and a body-to-inertial quaternion:
\[
\dot{\mathbf{p}}_n = R_{b}^{n}(\mathbf{q})\mathbf{v}_b,\qquad
\dot{\mathbf{v}}_b = \frac{\mathbf{F}_b}{m} + R_{n}^{b}(\mathbf{q})\begin{bmatrix}0&0&g\end{bmatrix}^T - \boldsymbol{\omega}_b\times\mathbf{v}_b,
\]
\[
\dot{\mathbf{q}} = \frac{1}{2}\Omega(\boldsymbol{\omega}_b)\mathbf{q},\qquad
J\dot{\boldsymbol{\omega}}_b = \mathbf{M}_b - \boldsymbol{\omega}_b\times J\boldsymbol{\omega}_b.
\]
The simulator renormalizes the quaternion after every RK4 step.

The truth model computes airdata from body velocity,
\[
V=\|\mathbf{v}_b\|,\qquad \alpha=\tan^{-1}(w/u),\qquad \beta=\sin^{-1}(v/V),
\]
and dynamic pressure $\bar{q}=\frac{1}{2}\rho V^2$. Attached-flow lift, drag, pitching moment, side-force, roll-moment, and yaw-moment coefficients are evaluated from $\alpha,\beta$, nondimensional rates, and control deflections. A smooth stall gate
\[
\sigma(\alpha)=\left(1+\exp\left[-\frac{|\alpha|-\alpha_{\mathrm{stall}}}{\Delta\alpha_{\mathrm{stall}}}\right]\right)^{-1}
\]
blends the attached lift curve into a finite post-stall lift model, adds post-stall drag rise, applies pitch-break moment, reduces control effectiveness, and adds nonlinear lateral-directional coupling. Body aerodynamic forces are then
\[
\begin{bmatrix}F_X&F_Y&F_Z\end{bmatrix}^T
= \bar{q}S\begin{bmatrix}C_X&C_Y&C_Z\end{bmatrix}^T,
\]
with $C_X=-C_D\cos\alpha+C_L\sin\alpha$ and $C_Z=-C_D\sin\alpha-C_L\cos\alpha$. Moments are scaled by $\bar{q}S[bC_l,cC_m,bC_n]^T$. Propeller thrust acts along body $x$ and contributes a small pitching moment. The dataset stores both the truth coefficients and an attached-flow nominal coefficient set; their difference is the hidden nonlinear aerodynamic residual.

Two observation assumptions are evaluated. In the direct-state case, the method receives noisy measurements of all 13 states during training and the initial measured state during validation. In the mocap-derived case, the method receives position and quaternion attitude. Body velocity is reconstructed by differentiating mocap position and rotating the result into the body frame, while angular rate is reconstructed by differentiating the quaternion and solving the local quaternion-rate equation. These reconstructed velocity and rate channels are intentionally noisier than the direct-state channel.

\IfFileExists{fig/generated_aircraft6dof_validation_score_comparison.svg}{%
\begin{figure}[hbt!]
    \centering
    \includesvg[width=0.95\linewidth]{generated_aircraft6dof_validation_score_comparison}
    \caption{6-DOF baseline method validation scores. Direct-state rows score full-state mean NRMSE. The mocap-output row scores position/quaternion output NRMSE because it does not predict body velocities or angular rates. Lower is better.}
    \label{fig:aircraft6dof_validation_score_comparison}
\end{figure}
}{}

\IfFileExists{fig/generated_aircraft6dof_train_time_accuracy_tradeoff.svg}{%
\begin{figure}[hbt!]
    \centering
    \includesvg[width=0.95\linewidth]{generated_aircraft6dof_train_time_accuracy_tradeoff}
    \caption{Training-time versus validation-error tradeoff for the 6-DOF benchmark. Methods with mean-state NRMSE above 1.0 are treated as failed rollouts and summarized in the callout so the useful region remains visible. The red dashed horizontal line gives the attached-flow nominal baseline, which has no fitted training solve. Marker area increases with validation rollout time.}
    \label{fig:aircraft6dof_train_time_accuracy_tradeoff}
\end{figure}
}{}

\IfFileExists{fig/generated_aircraft6dof_method_score_heatmap_direct.svg}{%
\begin{figure}[hbt!]
    \centering
    \includesvg[width=0.95\linewidth]{generated_aircraft6dof_method_score_heatmap_direct}
    \caption{Direct-state 6-DOF benchmark heatmap sorted by validation score. The score is full-state mean NRMSE over the nonlinear stall validation rollouts.}
    \label{fig:aircraft6dof_method_score_heatmap_direct}
\end{figure}
}{}

\IfFileExists{fig/generated_aircraft6dof_method_score_heatmap_mocap.svg}{%
\begin{figure}[hbt!]
    \centering
    \includesvg[width=0.95\linewidth]{generated_aircraft6dof_method_score_heatmap_mocap}
    \caption{Mocap-derived 6-DOF benchmark heatmap sorted by validation score. These rows expose the same pilot commands as the direct-state case but initialize and train from position/attitude-derived state estimates.}
    \label{fig:aircraft6dof_method_score_heatmap_mocap}
\end{figure}
}{}

\IfFileExists{fig/generated_aircraft6dof_validation_trajectory_overlay.svg}{%
\begin{figure}[hbt!]
    \centering
    \includesvg[width=0.95\linewidth]{generated_aircraft6dof_validation_trajectory_overlay}
    \caption{Representative 6-DOF validation rollout for the best implemented direct and mocap-derived state-predictor baselines. Rollouts use only the validation initial condition and pilot-command sequence after training; the panels include position, speed, angle of attack, stall-gate activation, and body rates.}
    \label{fig:aircraft6dof_validation_trajectory_overlay}
\end{figure}
}{}

Website-ready benchmark data are exported by
\begin{verbatim}
./results.py web-data
\end{verbatim}
into \texttt{site/public/data/}. The static site can be served locally with
\begin{verbatim}
./results.py serve-site
\end{verbatim}
The method-contribution workflow is intentionally two-phase: public pull requests add method code and pass plugin smoke checks, while full benchmark results are regenerated later in a trusted maintainer commit.

\section{Problem Formulation: 3 DOF Aircraft Longitudinal Model}\label{sec:prob}
The different system identification approaches would be investigated using a 3 DOF longitudinal aircraft motion, which means the aircraft is in a wings-level condition (\cite{StevensLewis2015}). In such a state, roll angle $\phi = 0$, and $\dot{\theta}=Q$. The side slip is assumed to be small, hence $\gamma = \theta - \alpha$.\\
Let $\{\mathbf{e}\}$, $\{\mathbf{w}\}$, and $\{\mathbf{b}\}$ denote the inertial, wind, and body frames, respectively.
\subsection*{States and Identified Parameter Vector}
The four flight-dynamic states and actuated control inputs are
\begin{align}
    \mathbf{x}_f &= \begin{bmatrix} V & \alpha & \gamma & Q \end{bmatrix}^{\!\top}, &
    \mathbf{u}   &= \begin{bmatrix} T_{\mathrm{}} & \delta_{e} \end{bmatrix}^{\!\top},
\end{align}
where $T_{\mathrm{}}$ and $\delta_{e,\mathrm{}}$ are the post-actuator (realized)
thrust and elevator deflection.
The aerodynamic parameter vector to be identified is
\begin{equation}
    \boldsymbol{\theta}_p =
    \begin{bmatrix}
        C_{L_0} & C_{L_\alpha} & C_{D_0} & k & C_{M_0} & C_{M_\alpha} & C_{M_Q} & C_{M_e}
    \end{bmatrix}^{\!\top}.
\end{equation}
\subsection*{Aerodynamic Model and Loads}
\subsection*{Aerodynamic Model}

Dynamic pressure is $\bar{q} = \tfrac{1}{2}\rho V^{2}$.
Aerodynamic forces and the pitching moment are computed as
\begin{equation}
    L = C_{L}\,\bar{q}\,S,
    \qquad
    D = C_{D}\,\bar{q}\,S,
    \qquad
    M = C_{M}\,\bar{q}\,S,
    \label{eq:aero_loads}
\end{equation}
with conventional coefficient expansions
\begin{align}
    C_{L} &= C_{L_0} + C_{L_\alpha}\,\alpha,
    \label{eq:CL}\\
    C_{D} &= C_{D_0} + k\,C_{L}^{2},
    \label{eq:CD}\\
    C_{M} &= C_{M_0}
           + C_{M_\alpha}\,\alpha
           + C_{M_{\delta_e}}\,\delta_{e,\mathrm{act}}
           + C_{M_Q}\,Q.
    \label{eq:CM}
\end{align}

The nominal parameter vector above is the structure available to the classical
fixed-form methods. It is an attached-flow model: lift is linear in angle of
attack, drag is parabolic in lift, and pitching moment is affine in
$\alpha$, $Q$, and elevator. The true simulator intentionally contains
additional unknown nonlinear aerodynamic terms,
\[
\mathbf{C}_{\mathrm{true}}(\mathbf{x},\mathbf{u})
=
\mathbf{C}_{\mathrm{nom}}(\mathbf{x},\mathbf{u};\boldsymbol{\theta}_p)
+\Delta\mathbf{C}_{\mathrm{aero}}(\mathbf{x},\mathbf{u}),
\]
where $\Delta\mathbf{C}_{\mathrm{aero}}=[\Delta C_L,\Delta C_D,\Delta C_M]^T$ is
not supplied to the identification algorithms. The hidden residual includes a
smooth stall model. Define
\[
\alpha_{\mathrm{stall}}=12.0^\circ,\qquad
w_{\mathrm{stall}}=2.0^\circ,\qquad
C_{L,\max}=1.05,\qquad
s_\alpha=\tanh(\alpha/0.015),
\]
\[
g_s =
\left[
1+\exp\left(-\frac{|\alpha|-\alpha_{\mathrm{stall}}}{w_{\mathrm{stall}}}\right)
\right]^{-1},
\qquad
\epsilon_\alpha =
w_{\mathrm{stall}}
\log\!\left(1+\exp\left(\frac{|\alpha|-\alpha_{\mathrm{stall}}}{w_{\mathrm{stall}}}\right)\right),
\]
where the code clips the exponential arguments for numerical stability. The
true lift coefficient is a smooth attached/separated-flow blend,
\begin{align}
C_{L,\mathrm{sep}} &=
s_\alpha(C_{L,\max}-1.35\epsilon_\alpha)
+0.035\tanh(1.5\delta_r)
+0.018\tanh(q_r),\\
C_{L,\mathrm{true}} &=
(1-g_s)C_{L,\mathrm{nom}}
+g_s C_{L,\mathrm{sep}}
+ (1-0.65g_s)\left[0.010\sin(1.8\delta_r)+0.010\alpha_r q_r\right],
\end{align}
and $\Delta C_L=C_{L,\mathrm{true}}-C_{L,\mathrm{nom}}$. Drag and pitching
moment use the residual forms
\begin{align}
\Delta C_D &=
0.0035\delta_r^2
+0.0025\alpha_r^4
+g_s\left(0.090+0.95\epsilon_\alpha+4.50\epsilon_\alpha^2\right)
+0.016g_s\max(-V_r,0),\\
\Delta C_M &=
-0.006\alpha_r^3
+0.006\alpha_r\delta_r
+0.004\tanh(q_r)|\alpha_r|
+0.0025\sin(2\gamma)
-g_s\left(0.075+0.55\epsilon_\alpha\right)s_\alpha
-0.018g_s\tanh(\delta_r),
\end{align}
with normalized variables
\[
V_r=(V-15)/15,\qquad
\alpha_r=\alpha/0.20,\qquad
q_r=\mathrm{sat}(Q/0.8),\qquad
\delta_r=\delta_e/0.25.
\]
This structure is intended to approximate a small RC trainer rather than a
transport-aircraft polar: the attached-flow lift slope is retained near trim,
lift approaches a finite separated-flow plateau near $C_L\approx 1$, drag rises
rapidly in the separated regime, and the pitching moment develops a nose-down
break for positive high-$\alpha$ flight. The terms create model-form error for
fixed-structure OEM and equation-error fits, while providing a known ground
truth for evaluating coefficient closure and residual-learning methods. The
SAFE-on aggressive case is expected to suppress part of this excitation because
the hidden controller commands recovery before the aircraft spends much time in
the high-$\alpha$ regime.

\paragraph{Note on pitching-moment normalisation.}
The expression $M = C_{M}\bar{q}S$ in \eqref{eq:aero_loads} omits the
reference chord $\bar{c}$ that appears in the standard convention ($M_{\mathrm{std}} = C_{M}^{\mathrm{std}}\bar{q}S\bar{c}$).
The simulation does not define $\bar{c}$ as a separate parameter; all
pitching-moment derivatives in \eqref{eq:CM} therefore implicitly absorb one
power of $\bar{c}$ and their numerical values are not directly comparable to
standard tabulated stability derivatives.
In particular, $C_{M_Q}$ as used here is a {dimensional} pitch-damping
derivative with units $\mathrm{s\,rad^{-1}}$, absorbing the factor
$\bar{c}/(2V)$ from the conventional nondimensional pitch rate
$\hat{q} = Q\bar{c}/(2V)$.
\subsection*{Equations of Motion (3-DoF Longitudinal)}
The equations governing the pure longitudinal motion are:
\begin{equation}
\boxed{
\begin{aligned}
\dot{V}      &= \frac{-D + T\cos\alpha - mg\sin\gamma}{m},\\[3pt]
\dot{\gamma} &= \frac{L + T\sin\alpha - mg\cos\gamma}{mV},\\[3pt]
\dot{Q}      &= \frac{M}{J_y},\\[3pt]
\dot{\alpha} &= Q - \dot{\gamma}.
\end{aligned}}
\label{eq:eom}
\end{equation}
\subsection*{Simulation Constants and Defaults}
Table~\ref{tab:simulation_defaults} lists the default constants used by the
current Python benchmark. These values are not intended to represent a specific
certified aircraft; they define a reproducible small fixed-wing longitudinal
test problem against which the identification methods are compared.

\begingroup
\small
\begin{longtable}{@{}p{0.24\linewidth}p{0.33\linewidth}p{0.35\linewidth}@{}}
\caption{Default simulator and benchmark constants used in the generated datasets.}
\label{tab:simulation_defaults}\\
\toprule
Quantity & Value & Code source or role \\
\midrule
\endfirsthead
\toprule
Quantity & Value & Code source or role \\
\midrule
\endhead
Mass $m$ & $1.0$ kg & \texttt{Aircraft.mass} \\
Pitch inertia $J_y$ & $0.15$ kg m$^2$ & \texttt{Aircraft.jy} \\
Reference area $S$ & $0.25$ m$^2$ & \texttt{Aircraft.wing\_area} \\
Air density $\rho$ & $1.225$ kg/m$^3$ & \texttt{Aircraft.rho} \\
Gravity $g$ & $9.81$ m/s$^2$ & \texttt{Aircraft.gravity} \\
Nominal $C_{L0}, C_{L\alpha}$ & $0.10,\;3.00$ & \texttt{NominalAero} \\
Nominal $C_{D0}, k$ & $0.030,\;0.10$ & \texttt{NominalAero} \\
Nominal $C_{M0}, C_{M\alpha}$ & $0.010,\;-0.100$ & \texttt{NominalAero} \\
Nominal $C_{MQ}, C_{M\delta_e}$ & $-0.100,\;0.100$ & \texttt{NominalAero} \\
Sample time $\Delta t$ & $0.01$ s & locked 100 Hz mocap rate \\
Trajectory duration & $40$ s & \texttt{SimulationConfig.duration} \\
Training/validation trials & $64/16$ & \texttt{SimulationConfig} \\
Direct-state noise std. & $[0.08,0.0035,0.0035,0.012]$ & $[V,\alpha,\gamma,Q]$ \\
Mocap noise std. & $0.002$ m, $0.0015$ rad & $[p_x,p_z,\theta]$ \\
Mocap smoothing window & $21$ samples & moving average in simulator; Savitzky--Golay in comparison suite \\
Thrust actuator & $\tau_T=0.12$ s, $|\dot{T}|\le 7.0$ N/s, $T\in[0,3]$ N & actuator lag/rate/saturation \\
Elevator actuator & $\tau_e=0.055$ s, $|\dot{\delta}_e|\le 2.4$ rad/s, $\delta_e\in[-0.35,0.35]$ rad & actuator lag/rate/saturation \\
SAFE blend $\beta$ & $0.85$ & hidden-controller simulator \\
SAFE pitch gain $K_\theta$ & $3.0$ & hidden-controller simulator \\
SAFE elevator gain $K_q$ & $0.105$ & hidden-controller simulator \\
AS3X pitch-rate damping & $0.018$ & hidden-controller simulator \\
Recovery gates & $\alpha_{\lim}=0.110$ rad, $V_{\lim}=13.0$ m/s & hidden-controller simulator \\
\bottomrule
\end{longtable}
\endgroup
\subsection*{Indoor Disturbance Assumption}

The motivating experiment takes place in an indoor aircraft hangar, so
atmospheric gusts are not treated as a primary uncertainty source in the
current benchmark. The generated datasets therefore disable the colored
process-disturbance channel by default. The dominant practical mismatch is
instead input uncertainty: the experimenter observes the external pilot
command, while the off-the-shelf vehicle may internally reshape that command
through actuator dynamics and a proprietary SAFE/AS3X stabilization loop.

\subsection*{Measurement and Input Model}

The oracle diagnostic channel provides the four flight-dynamic states
corrupted by additive white Gaussian noise:
\begin{equation}
    \mathbf{y} =
    \begin{bmatrix} V \\ \alpha \\ \gamma \\ Q \end{bmatrix}
    + \boldsymbol{v},
    \qquad
    \boldsymbol{v} \sim
    \mathcal{N}\!\left(\mathbf{0},\;
        R\right),
    \quad
    R = \mathrm{diag}\!\left(
        \sigma_{V}^{2},\;
        \sigma_{\alpha}^{2},\;
        \sigma_{\gamma}^{2},\;
        \sigma_{Q}^{2}
    \right),
    \label{eq:meas}
\end{equation}
where the direct-state channel is used only to diagnose algorithm behavior.
The practical input supplied to the identification methods is the external
pilot command,
\[
\mathbf{u}_{\mathrm{cmd}} = [T_{\mathrm{cmd}},\delta_{e,\mathrm{pilot}}]^T.
\]
When SAFE is disabled, this command passes only through actuator lag and rate
limits before entering the plant. When SAFE is enabled, the simulator inserts
a hidden longitudinal projection of the proprietary controller between the
pilot command and the elevator actuator. The true realized input
$\mathbf{u}_{\mathrm{act}}$ and the hidden SAFE correction are retained in the
dataset for scoring and diagnostics, but the practical identification runs use
only $\mathbf{u}_{\mathrm{cmd}}$.

\subsection*{Hidden SAFE/AS3X Pitch Loop}
The real Sport Cub S2 firmware is proprietary, so the benchmark uses an
engineering approximation of the longitudinal part of a SAFE/AS3X controller.
This model is not intended to claim the true Horizon/Spektrum implementation.
It creates a controlled hidden-input problem that is plausible for an
off-the-shelf trainer: the pilot commands throttle and elevator, while an
internal pitch-attitude/rate loop reshapes elevator before the plant sees it.

The estimated pitch attitude is
\[
\theta = \alpha + \gamma,
\]
and throttle creates a mode-dependent pitch bias,
\begin{equation}
    \theta_{\mathrm{thr}} =
    \frac{1}{2}\,\mathrm{interp}
    \left(T_{\mathrm{cmd}}/T_{\max};
    [0,0.5,1],[-5^\circ,0^\circ,8^\circ]\right).
\end{equation}
The pilot elevator command is mapped to a bounded pitch command,
\begin{equation}
    \theta_{\mathrm{cmd}}
    = \theta_0 + \theta_{\mathrm{thr}}
    + K_{\delta}\,\delta_{e,\mathrm{stick}},
    \qquad
    \delta_{e,\mathrm{stick}}
    = \mathrm{sat}\left(\frac{\delta_{e,\mathrm{pilot}}-\delta_{e,\mathrm{trim}}}{0.20}\right).
\end{equation}
The SAFE outer loop produces a pitch-rate command and elevator correction,
\begin{align}
    q_{\mathrm{cmd}} &= \mathrm{sat}\left(K_\theta(\theta_{\mathrm{cmd}}-\theta)\right),\\
    \delta_{e,\mathrm{SAFE}} &= \delta_{e,\mathrm{trim}} + K_q(q_{\mathrm{cmd}}-q).
\end{align}
The AS3X-like inner-loop contribution is modeled as direct pilot elevator plus
pitch-rate damping,
\begin{equation}
    \delta_{e,\mathrm{AS3X}} = \delta_{e,\mathrm{pilot}} - K_{\mathrm{AS3X}}q.
\end{equation}
Finally, a smooth recovery gate activates a stronger pitch-recovery law near
high angle of attack or low airspeed,
\begin{equation}
    r = 1 - (1-\sigma_\alpha)(1-\sigma_V),
    \quad
    \sigma_\alpha=\sigma\!\left(\frac{\alpha-\alpha_{\lim}}{s_\alpha}\right),
    \quad
    \sigma_V=\sigma\!\left(\frac{V_{\lim}-V}{s_V}\right),
\end{equation}
and the internal elevator command becomes
\begin{equation}
    \delta_{e,\mathrm{int}}
    = (1-r)\left(\beta\delta_{e,\mathrm{SAFE}}+(1-\beta)\delta_{e,\mathrm{AS3X}}\right)
    + r\delta_{e,\mathrm{panic}}.
    \label{eq:safe_loop}
\end{equation}
The internal throttle is not closed-loop controlled in this benchmark,
$T_{\mathrm{int}}=T_{\mathrm{cmd}}$. This matches the intended use of the
3-DoF model: we test whether unobserved longitudinal elevator reshaping can
confound or be learned by system-identification methods, while avoiding an
unjustified full 6-DoF roll/yaw SAFE model.

\subsection*{Actuator Dynamics}
Both effectors are modelled as first-order lags with position saturation and rate
limiting:
\begin{align}
    \dot{T}_{\mathrm{act}}
        &= \frac{T_{\mathrm{cmd}} - T_{\mathrm{act}}}{\tau_T},
    & T_{\mathrm{act}}   &\in [0,\;3]\;\text{N},
    & |\dot{T}_{\mathrm{act}}|   &\le 7.0\;\text{N/s},
    & \tau_T   &= 0.12\;\text{s},
    \label{eq:Tact}\\[4pt]
    \dot{\delta}_{e,\mathrm{act}}
        &= \frac{\delta_{e,\mathrm{cmd}} - \delta_{e,\mathrm{act}}}{\tau_e},
    & \delta_{e,\mathrm{act}} &\in [-0.35,\;0.35]\;\text{rad},
    & |\dot{\delta}_{e,\mathrm{act}}| &\le 2.4\;\text{rad/s},
    & \tau_e   &= 0.055\;\text{s}.
    \label{eq:eact}
\end{align}
Saturation and rate limits introduce dynamics absent from the nominal
model~\eqref{eq:eom}, causing structural mismatch when identification uses only
the external pilot command.
\subsection*{Test Cases} \label{subsec:testcase}
The current comparison uses five reproducible datasets. The \texttt{open\_loop}
dataset uses randomized pilot commands with SAFE disabled. The
\texttt{sine\_sweep} dataset uses elevator sine sweeps with SAFE disabled for
frequency-response diagnostics. The \texttt{open\_loop\_safe} and
\texttt{sine\_sweep\_safe} datasets repeat those excitation families with the
hidden SAFE/AS3X pitch loop enabled. The \texttt{safe\_loop} dataset uses
recovery-probe pilot inputs with pull-up and throttle-reduction segments that
make the hidden loop observable. These cases separate three effects: excitation
quality, motion-capture state reconstruction, and unobserved proprietary
command reshaping.

For the motion-capture experiments motivating this benchmark, the primary measured outputs are not the aerodynamic states themselves. In the PURT pylon-racing and indoor testbed setting \citep{sribunma2026indoor}, the aircraft are small commercial fixed-wing vehicles and the only reliable ground-truth measurement is the external motion-capture track. No pitot-static system, onboard IMU, or direct angle-of-attack sensor is assumed. A typical laboratory system provides inertial position and attitude at approximately 100 Hz,
\[
\mathbf{y}_{\mathrm{mocap},k} =
\begin{bmatrix}
p_{x,k} & p_{z,k} & \theta_k
\end{bmatrix}^{T}
\;+\;
\boldsymbol{\eta}_k,
\]
while the longitudinal states $\mathbf{x}=[V,\alpha,\gamma,q]^T$ are latent variables. These states can be reconstructed approximately from smoothed finite differences of position and attitude, but this operation amplifies measurement noise and introduces smoothing-dependent phase lag and bandwidth loss. Therefore, the numerical study separates three cases: an oracle state-measurement case using simulated noisy $V,\alpha,\gamma,q$ directly, a derived-state case using finite-difference estimates from mocap data, and a preferred mocap-output case in which the model propagates latent flight states and compares the predicted output $h(\mathbf{x})=[p_x,p_z,\theta]^T$ directly against motion-capture measurements. The simulator retains the true state and aerodynamic residuals only for scoring and diagnostic plots; these channels are not assumed to be available to the identification algorithms in the mocap-output comparison.

This distinction becomes more important as the observation rate is reduced. At 100 Hz the Nyquist frequency is 50 Hz, which is typically well above the dominant rigid-body longitudinal content of a small fixed-wing aircraft. At 10 Hz, representative of a slow GPS-like position stream, the Nyquist frequency falls to 5 Hz and derivative estimates become much more sensitive to filtering choices. At 2 Hz, the Nyquist frequency is only 1 Hz, so short-period or actuator-driven dynamics can be attenuated, aliased, or rendered practically unobservable from position and attitude alone. The lower-rate cases are therefore treated as observation-quality studies rather than changes to the aircraft dynamics.

\subsection*{Benchmark Protocol}
The current benchmark is organized around five complementary questions. First, under idealized state availability, how well do OEM, equation-error least squares, SINDy, PINN-style coefficient closure, UDE residual learning, and neural coefficient surrogates handle nominal-plus-hidden nonlinear aerodynamic terms? This oracle comparison isolates model-form error from state-reconstruction error. Second, under realistic motion-capture observation, how much performance is lost when the same methods are forced to use derived states, and how much can be recovered by fitting the measured position/attitude outputs directly? Third, when a hidden SAFE/AS3X-like loop reshapes elevator commands, can the methods distinguish missing aerodynamics from missing input dynamics? Fourth, when the maneuver deliberately drives larger angle-of-attack, pitch-rate, elevator, and speed excursions, do methods that look adequate on local trajectories still recover the nonlinear aerodynamic residuals? Fifth, when the training set spans a deliberately wider initial-condition distribution, does one global model remain adequate or does a stitched collection of local models improve extrapolation across operating regimes?

For this reason the benchmark reports separate rows for direct-state and mocap-derived-state training where applicable. The direct-state case uses noisy measurements of $V,\alpha,\gamma,q$ and should be interpreted as an upper-bound diagnostic for algorithms that require state histories or state derivatives. The mocap-derived case uses $p_x,p_z,\theta$ sampled at 100 Hz and reconstructs $V,\alpha,\gamma,q$ by smoothing and finite differencing; this case is intentionally included because many practical pipelines use this preprocessing step, even though it is numerically fragile. The mocap-output OEM case instead keeps $V,\alpha,\gamma,q$ as latent states and optimizes the output error in $p_x,p_z,\theta$ using a CasADi/IPOPT nonlinear program. This is the preferred deterministic baseline for laboratory motion-capture data because it avoids differentiating the measured position and attitude channels.

The observation-rate study decimates the same 100 Hz mocap stream to 10 Hz and 2 Hz before applying the state-reconstruction pipeline. This does not change the aircraft trajectory or aerodynamic model; it only changes the information available to the estimator. The resulting degradation quantifies the practical observability and sampling limitations imposed by slower sensors. As the observation rate approaches the dominant aircraft modes, Nyquist limitations and filtering tradeoffs become central rather than incidental implementation details.

The method comparison is now organized as a train-once, validate-everywhere
benchmark. Each method is assigned one identification dataset that matches its
intended experimental design, and the fitted model is then rolled out on every
validation condition. Local linear, frequency-domain, subspace, and
model-stitching baselines train on the trim-grid dataset because those methods
are normally supported by small-deviation excitation near selected operating
points. Nonlinear aerodynamic residual and coefficient-closure methods train on
the aggressive SAFE-off dataset because they need high-angle-of-attack,
pitch-rate, speed, and control variation to expose the hidden aerodynamic
terms. Hidden-control methods train on the SAFE recovery-probe dataset because
that is the dataset where the external pilot command differs systematically
from the internal elevator command. The nominal model has no fitted training
solve. This protocol prevents the easier in-distribution experiment in which a
separate model is trained for every validation maneuver.

The local open-loop dataset intentionally remains a near-trim maneuver. It is
included because small-amplitude identification is still a common aircraft
SysID use case and because it provides a sanity check where global linear
models should be competitive. The aggressive sine-sweep and aggressive
nonlinear datasets use a stratified initial-condition sampler instead of a
single Gaussian perturbation about trim. SAFE-off aggressive trials are drawn
from local-trim, low-speed/high-$\alpha$, pull-up, recovery, and fast-dive
bands: approximately $V\in[7.5,29]$ m/s,
$\alpha\in[-5^\circ,20^\circ]$, $\gamma\in[-35^\circ,32^\circ]$, and
$q\in[-95,90]^\circ$/s over the union of regimes. The validation split shifts
some high-$\alpha$ bands upward to test interpolation and limited
extrapolation. SAFE-on trials use a broad but recoverable subset of the same
idea so that the hidden recovery loop can visibly suppress high-$\alpha$
excursions. The trim-grid dataset is different from the generic open-loop
dataset: it consists of small multisine and doublet perturbations about several
trim-like operating points and is used to support the local-model and
frequency-domain training protocols. This design separates the question ``does
a local linear model work near trim?'' from the harder question ``does one
identified model remain adequate across trim, stall onset, recovery, and speed
variation?''

\subsection*{Implemented Methods}
Each method is run from its assigned training split and then evaluated on the
common validation suite. Unless stated otherwise, methods are given the external
pilot command $\mathbf{u}_{\mathrm{cmd}}$
and either direct noisy states or states reconstructed from 100 Hz motion capture.
The validation score is always computed against the simulator truth and is not a
training objective. It is the mean normalized state RMSE over the validation
rollout,
\[
J_{\mathrm{val}}
=
\frac{1}{4}\sum_{j=1}^{4}
\frac{
\sqrt{N^{-1}\sum_{i,k}
\left(\hat{x}_{i,k,j}-x_{i,k,j}^{\mathrm{true}}\right)^2}
}{
\max_{i,k} x_{i,k,j}^{\mathrm{true}}
-\min_{i,k} x_{i,k,j}^{\mathrm{true}}
},
\]
with the denominator replaced by one only for numerically constant channels.
Thus the score is dimensionless: it is not the RMSE of one state, but the
average of the range-normalized RMSEs for $V,\alpha,\gamma$, and $q$. Primary
ranking figures and tables report open-loop model rollouts driven only by the
supplied pilot-command history after initialization. Methods may use the
validation initial condition. Lagged input-output methods initialize their
required state history from that same initial condition and the pilot-command
history rather than from additional validation measurements, so validation
measurements are not fed back into the model after rollout begins.

\paragraph{Nominal model.}
The nominal baseline propagates \eqref{eq:eom} using the nominal aerodynamic
coefficients and no residual correction. It is included to expose the size of the
combined aerodynamic and hidden-input model error.

\paragraph{Equation-error least squares.}
The equation-error method smooths the measured state history with a
Savitzky--Golay filter, differentiates it, inverts the equations of motion to
obtain inferred $C_L$, $C_D$, and $C_M$, and solves linear least-squares problems
for the nominal coefficient vector. This is fast and interpretable, but it is
biased by measurement noise, derivative noise, hidden actuator dynamics, and
SAFE command reshaping.

\paragraph{Output-error methods.}
The single-shooting OEM fits the nominal aerodynamic parameters by minimizing
trajectory output error under RK4-integrated dynamics using CasADi/IPOPT. The
mocap-output OEM keeps $V,\alpha,\gamma,q$ latent and compares the integrated
outputs $p_x,p_z,\theta$ directly to measured motion-capture outputs. The
\texttt{OEM-HiddenController} variant additionally estimates the parameters of
the SAFE-like pitch loop in \eqref{eq:safe_loop}. This last variant is an
oracle-structure benchmark: in a real proprietary aircraft, the controller form
is not known. It should therefore be read as an upper bound for a correctly
structured hidden-controller estimator, not as a fair black-box discovery method.

\paragraph{EKF parameter identification.}
The EKF row initializes the nominal aerodynamic parameter vector from equation
error and then runs an augmented-state EKF with a CasADi
automatic-differentiated RK4 transition and parameter Jacobian on the training
data. The augmented filter estimates the aerodynamic parameter vector from
training measurements only. After training, the parameter vector is frozen and
validation uses the same pilot-command-only open-loop rollout protocol as the
other dynamics-model rows.

\paragraph{Frequency-Welch and Frequency-Stitching.}
The frequency-domain baseline estimates averaged input/state spectra using
Welch-style windows, discards low-coherence frequency bins, and fits a local
continuous-time linear model from the weighted complex relation
$j\omega X(\omega)=AX(\omega)+BU(\omega)$. This is not CIFER; it is a compact
spectral baseline with some of the same diagnostic ingredients. Its result
should not be interpreted as a judgement of CIFER performance. The
Frequency-Stitching variant uses the same trim-grid training set as the local
linear stitching baseline, fits one Welch linearization per trim-grid group,
and blends the local derivative predictions with the same smooth scheduling
weights. It is included to test whether the poor global frequency-domain row is
mostly a single-LTI-model limitation or a deeper consequence of limited
coherence, derivative noise, and non-frequency-optimized excitation.

\paragraph{Subspace, Koopman, local linear, and stitched models.}
The linear state-space baseline fits a one-step affine discrete model by least
squares. The Hankel/SVD subspace baseline constructs lagged input/output blocks,
projects them onto a low-dimensional latent basis, and fits linear latent
dynamics. The Koopman-EDMD baseline uses the same nonlinear candidate functions
as SINDy, but fits a lifted discrete-time one-step map rather than sparse
continuous-time derivatives. The model-stitching baseline clusters training
trajectories by initial condition, fits multiple local affine state-space
models, and blends their one-step predictions with smooth scheduling weights
during validation. These methods are included because they can be
strong predictors even when they do not recover physically meaningful
aerodynamic parameters.

\paragraph{SINDy and symbolic stepwise regression.}
The SINDy row uses a structured sparse coefficient formulation. It regresses the
EOM-inferred aerodynamic coefficients $C_L$, $C_D$, and $C_m$ on protected
low-order aircraft terms plus a sparse residual library in one fit. Protected
terms are always retained, while residual-library terms are selected by
sequential thresholded least squares. The stepwise symbolic baseline remains a
black-box sparse derivative model. These sparse-library methods are sensitive
to derivative quality, excitation, and to whether the library contains the true
hidden mechanism. Saturations, rate limits, and recovery gates are especially
hard because they appear only near
boundaries unless the experiment deliberately visits those regions.

\paragraph{PINN, UDE, GP/RBF, and neural surrogates.}
The coefficient-closure PINN and UDE residual models keep the longitudinal
equations of motion and learn missing aerodynamic or state-derivative terms with
small neural networks trained on GPU. The GP/RBF coefficient closure is a
sparse kernel surrogate trained on the same equation-of-motion-inferred
coefficient residuals as the PINN-style closure. The hidden-control UDE instead learns an
effective throttle/elevator correction from equations-of-motion residuals, and
the hidden-elevator PINN constrains that correction to the elevator channel using
the pitch-moment equation. The supervised coefficient surrogate uses synthetic
labels for $\Delta\mathbf{C}_{\mathrm{aero}}$ and is therefore an upper-bound
diagnostic rather than a flight-data-only method.

\subsection*{Code-Matched Benchmark Specification}
This subsection is intended to make the paper auditable against the Python
implementation. The symbol $\mathbf{z}_k=[\mathbf{x}_k^T,\mathbf{u}_k^T]^T$
denotes the state-input sample used by the regression and neural methods. When
the command-line option \texttt{--input-channel u\_cmd} is used, the benchmark
intentionally replaces the internal input array passed to the methods with the
external pilot command. Thus the code variable named \texttt{u\_act} inside
\texttt{comparison\_suite.py} should be read as the input channel supplied to
the estimator, not necessarily the true actuator input. This convention is what
creates the practical hidden-controller problem in the SAFE-on cases.

The common smoothing and derivative pipeline used by derivative-based methods is
\begin{align}
    \widetilde{\mathbf{x}}_{1:N}
        &= \mathrm{SG}_{w,p}\!\left(\mathbf{y}_{1:N}\right),\\
    \dot{\widetilde{\mathbf{x}}}_{k}
        &= \mathrm{gradient}\!\left(\widetilde{\mathbf{x}}_{1:N},\Delta t\right)_k,
        \label{eq:code_derivative_pipeline}
\end{align}
where the current code uses a Savitzky--Golay window $w=21$ and polynomial
order $p=3$ unless overridden. For mocap-derived runs, $\mathbf{y}_{1:N}$ is
first replaced by the reconstructed state history obtained from
$p_x,p_z,\theta$.

\paragraph{Nominal rollout.}
The nominal method has no training step. It sets
$\hat{\boldsymbol{\theta}}=\boldsymbol{\theta}_{\mathrm{nom}}$ and propagates
\[
    \hat{\mathbf{x}}_{k+1}
    =\Phi_{\mathrm{RK4}}
    \left(\hat{\mathbf{x}}_k,\mathbf{u}_k,\mathbf{u}_{k+1};
    \boldsymbol{\theta}_{\mathrm{nom}}\right),
    \qquad
    \hat{\mathbf{x}}_0=\mathbf{y}_0.
    \label{eq:code_nominal}
\]
This is implemented by \texttt{run\_nominal}, \texttt{theta\_dynamics}, and
\texttt{rk4\_step}.

\paragraph{Equation-error least squares.}
The equation-error method first algebraically infers aerodynamic coefficients
from smoothed state derivatives:
\begin{align}
    \widehat{D}_k
        &= T_k\cos\widetilde{\alpha}_k
         -mg\sin\widetilde{\gamma}_k
         -m\dot{\widetilde{V}}_k,\\
    \widehat{L}_k
        &= m\widetilde{V}_k\dot{\widetilde{\gamma}}_k
         -T_k\sin\widetilde{\alpha}_k
         +mg\cos\widetilde{\gamma}_k,\\
    \widehat{M}_k
        &= J_y\dot{\widetilde{Q}}_k,\\
    \widehat{\mathbf{C}}_k
        &= [\widehat{L}_k,\widehat{D}_k,\widehat{M}_k]^T
           /(\bar{q}_k S).
        \label{eq:code_inferred_coefficients}
\end{align}
It then solves three ordinary least-squares problems,
\begin{align}
    [\widehat{C}_{L0},\widehat{C}_{L\alpha}]^T
        &= \arg\min_{\beta}\sum_k
        \left(\widehat{C}_{L,k}-[1,\widetilde{\alpha}_k]\beta\right)^2,\\
    [\widehat{C}_{D0},\widehat{k}]^T
        &= \arg\min_{\beta}\sum_k
        \left(\widehat{C}_{D,k}-[1,\widehat{C}_{L,k}^{\,2}]\beta\right)^2,\\
    [\widehat{C}_{M0},\widehat{C}_{M\alpha},
     \widehat{C}_{MQ},\widehat{C}_{M\delta_e}]^T
        &= \arg\min_{\beta}\sum_k
        \left(\widehat{C}_{M,k}
        -[1,\widetilde{\alpha}_k,\widetilde{Q}_k,\delta_{e,k}]\beta\right)^2.
        \label{eq:code_equation_error}
\end{align}
The fitted vector is clipped to the code bounds
$[-0.5,0,0,0,-0.2,-1,-1,-0.5]\le\theta\le
[0.5,6,0.2,0.5,0.2,0.5,0.1,0.5]$. This matches
\texttt{infer\_coefficients} and \texttt{fit\_equation\_error}.

\paragraph{Output-error single shooting.}
The coded single-shooting OEM estimates one global aerodynamic parameter vector
and one initial state for each selected training trajectory:
\begin{equation}
    \min_{\theta,\{\mathbf{x}_{0}^{(r)}\}}
    \sum_{r\in\mathcal{R}}\sum_{k\in\mathcal{K}}\sum_{j=1}^{4}
    \left(
    \frac{\hat{x}_{k,j}^{(r)}-y_{k,j}^{(r)}}{\sigma_j}
    \right)^2,
    \qquad
    \hat{\mathbf{x}}_{k+1}^{(r)}
    =
    \Phi_{\mathrm{RK4}}
    (\hat{\mathbf{x}}_{k}^{(r)},\mathbf{u}_{k}^{(r)},\mathbf{u}_{k+1}^{(r)};
    \theta).
    \label{eq:code_oem_ss}
\end{equation}
The direct-state noise weights are
$\sigma=[0.08,0.0035,0.0035,0.012]$. The current comparison uses CasADi/IPOPT,
bounded parameters, bounded initial states, a selected number of trajectories,
and a fit stride. This is the implementation in
\texttt{run\_oem\_casadi}.

\paragraph{Mocap-output OEM.}
The mocap-output OEM uses the same aerodynamic parameterization but changes the
output map. Position is integrated by trapezoidal quadrature,
\begin{align}
    \hat{p}_{x,k+1}
        &= \hat{p}_{x,k}
        +\frac{\Delta t}{2}
        (\hat{V}_k\cos\hat{\gamma}_k
        +\hat{V}_{k+1}\cos\hat{\gamma}_{k+1}),\\
    \hat{p}_{z,k+1}
        &= \hat{p}_{z,k}
        +\frac{\Delta t}{2}
        (\hat{V}_k\sin\hat{\gamma}_k
        +\hat{V}_{k+1}\sin\hat{\gamma}_{k+1}),\\
    \hat{\theta}_k &= \hat{\alpha}_k+\hat{\gamma}_k.
\end{align}
The optimization objective is
\begin{equation}
    \min_{\theta,\{\mathbf{x}_0^{(r)}\}}
    \sum_{r,k}
    \left[
    \left(\frac{\hat{p}_{x,k}^{(r)}-p_{x,k}^{(r)}}{\sigma_p}\right)^2
    +\left(\frac{\hat{p}_{z,k}^{(r)}-p_{z,k}^{(r)}}{\sigma_p}\right)^2
    +\left(\frac{\hat{\theta}_{k}^{(r)}-\theta_{k}^{(r)}}{\sigma_\theta}\right)^2
    \right],
    \label{eq:code_oem_mocap}
\end{equation}
with $\sigma_p=0.002$ m and $\sigma_\theta=0.0015$ rad by default. This is
\texttt{run\_oem\_mocap\_output\_casadi}.

\paragraph{Variational-Mocap sparse MAP surrogate.}
The coded variational row is a deterministic sparse-NLP analogue of a
variational smoother. It optimizes aerodynamic parameters, latent flight-state
nodes, and latent position nodes:
\begin{align}
    \min_{\theta,\{\mathbf{x}_k^{(r)},\mathbf{p}_k^{(r)}\}}
    \sum_{r,k}
    \left\|\frac{h(\mathbf{x}_k^{(r)},\mathbf{p}_k^{(r)})
    -\mathbf{y}_{\mathrm{mocap},k}^{(r)}}{\boldsymbol{\sigma}_y}\right\|_2^2
    &+
    \sum_{r,k}
    \left\|\frac{\mathbf{x}_{k+1}^{(r)}
    -\Phi_{\mathrm{RK4}}(\mathbf{x}_{k}^{(r)},\mathbf{u}_{k}^{(r)},\mathbf{u}_{k+1}^{(r)};\theta)}
    {\boldsymbol{\sigma}_x}\right\|_2^2 \nonumber\\
    &+
    \sum_{r,k}
    \left\|\frac{\mathbf{p}_{k+1}^{(r)}
    -\Psi_p(\mathbf{p}_{k}^{(r)},\mathbf{x}_{k}^{(r)},\mathbf{x}_{k+1}^{(r)})}
    {\sigma_{p,\mathrm{proc}}}\right\|_2^2 .
    \label{eq:code_variational_mocap}
\end{align}
Here $h=[p_x,p_z,\alpha+\gamma]^T$ and $\Psi_p$ is the same trapezoidal
position update used by mocap-output OEM. This matches
\texttt{run\_variational\_mocap\_output\_casadi}; it should not be read as a
full posterior VI implementation because no covariance or amortized posterior
is learned.

\paragraph{OEM with hidden SAFE-like controller.}
The hidden-controller OEM uses the same objective as
\eqref{eq:code_oem_ss}, but the plant input is
\[
    \mathbf{u}_{\mathrm{int},k}
    =
    g_{\mathrm{SAFE}}(\mathbf{x}_k,\mathbf{u}_{\mathrm{cmd},k};
    \psi,\mathbf{u}_{\mathrm{trim}}),
\]
where $g_{\mathrm{SAFE}}$ has the same algebraic structure as
\eqref{eq:safe_loop}. The decision vector is
$[\theta^T,\psi^T,\mathbf{x}_{0}^{(1)T},\ldots]^T$ and the coded objective is
\begin{equation}
    \min_{\theta,\psi,\{\mathbf{x}_0^{(r)}\}}
    \sum_{r,k,j}
    \left(\frac{\hat{x}_{k,j}^{(r)}-y_{k,j}^{(r)}}{\sigma_j}\right)^2
    +10^{-2}\|\psi-\psi_0\|_2^2.
    \label{eq:code_oem_hidden}
\end{equation}
The controller vector is
$\psi=[\theta_0,K_T,K_\delta,K_\theta,K_q,K_{\mathrm{AS3X}},
\alpha_{\lim},V_{\lim},\Delta\theta_{\mathrm{panic}},
K_{\theta,\mathrm{panic}},K_{q,\mathrm{panic}},\beta]^T$.
This is an intentionally optimistic oracle-structure estimator implemented in
\texttt{run\_oem\_hidden\_controller}.

\paragraph{EKF parameter identification.}
The EKF parameter-identification row first fits $\theta$ with equation-error
LS, then runs an augmented-state EKF on the training split. The augmented state
is $\mathbf{z}_k=[\mathbf{x}_k^T,\boldsymbol{\theta}^T]^T$, where
$\boldsymbol{\theta}$ is modeled as a random walk with small process noise. The
RK4 transition and both Jacobians are generated by CasADi automatic
differentiation:
\begin{align}
    \hat{\mathbf{x}}_{k+1|k}
        &=F(\hat{\mathbf{x}}_{k|k},\mathbf{u}_k,\mathbf{u}_{k+1},
            \hat{\boldsymbol{\theta}}_{k|k},\Delta t),\\
    A_k &= \frac{\partial F}{\partial \mathbf{x}},
    \qquad
    B_k = \frac{\partial F}{\partial \boldsymbol{\theta}}.
    \label{eq:code_ekf}
\end{align}
Training measurements update both state and parameter estimates inside the EKF
recursion. The final parameter estimate is then frozen, and validation is scored
by integrating the identified dynamics open loop from the validation initial
condition using only the pilot-command sequence. This is
\texttt{run\_filter\_error\_ekf}, reported as \texttt{EKF-ParamID}.

\paragraph{Linear-SS.}
The local affine discrete state-space baseline solves
\begin{equation}
    \widehat{C}
    =
    \arg\min_C
    \sum_k
    \left\|
    \mathbf{x}_{k+1}
    -
    [1,\mathbf{x}_k^T,\mathbf{u}_k^T]C
    \right\|_2^2,
    \qquad
    \hat{\mathbf{x}}_{k+1}
    =
    [1,\hat{\mathbf{x}}_k^T,\mathbf{u}_k^T]\widehat{C}.
    \label{eq:code_linear_ss}
\end{equation}
Validation is then an open-loop rollout from the validation initial condition
using only the supplied pilot-command history. The method therefore does not
use validation measurements as feedback, but it also does not estimate
aerodynamic derivatives or hidden-controller parameters. Low validation error
for this row should be interpreted as strong in-distribution discrete
prediction, not as physical interpretability. This is
\texttt{run\_linear\_state\_space}.

\paragraph{Model-Stitching.}
The stitched local-model row uses a different training protocol from the other
rows. Rather than fitting on the same maneuver that defines each validation
case, it always trains on the trim-grid dataset: a collection of small-deviation
multisine and doublet records collected about several trim-like operating
points. This reflects the usual gain-scheduled aircraft-identification workflow:
collect local excitation near several operating points, fit local deviation
models, and then stitch those local models for broader-envelope prediction.
With $M$ local models, k-means assigns each trim-grid trajectory to a regime
index $r_i\in\{1,\ldots,M\}$ using its local operating state and control
$\boldsymbol{\rho}_i=[\mathbf{x}_{i,\mathrm{trim}}^T,
\mathbf{u}_{i,\mathrm{trim}}^T]^T$. For each regime, the code fits an affine
deviation model,
\begin{equation}
    \widehat{C}^{(m)}
    =
    \arg\min_C
    \sum_{i:r_i=m}\sum_k
    \left\|
    (\mathbf{x}_{i,k+1}-\mathbf{x}^{(m)}_{\mathrm{trim}})
    -
    [1,(\mathbf{x}_{i,k}-\mathbf{x}^{(m)}_{\mathrm{trim}})^T,
       (\mathbf{u}_{i,k}-\mathbf{u}^{(m)}_{\mathrm{trim}})^T]C
    \right\|_2^2
    +\lambda\|C\|_F^2.
    \label{eq:code_model_stitching_local}
\end{equation}
Each local model is associated with the operating point
$\boldsymbol{\mu}^{(m)}=[(\mathbf{x}^{(m)}_{\mathrm{trim}})^T,
(\mathbf{u}^{(m)}_{\mathrm{trim}})^T]^T$. During rollout, the local
predictions
\[
    \hat{\mathbf{x}}_{k+1}^{(m)}
    =
    \mathbf{x}^{(m)}_{\mathrm{trim}}
    +
    [1,(\hat{\mathbf{x}}_k-\mathbf{x}^{(m)}_{\mathrm{trim}})^T,
       (\mathbf{u}_k-\mathbf{u}^{(m)}_{\mathrm{trim}})^T]\widehat{C}^{(m)}
\]
are blended with normalized radial-basis scheduling weights
\begin{align}
    w_m(\hat{\mathbf{x}}_k,\mathbf{u}_k)
    &=
    \frac{
    \exp\left[-\frac{1}{2\ell^2}
    \left\|S_\rho^{-1}([\hat{\mathbf{x}}_k^T,\mathbf{u}_k^T]^T-\boldsymbol{\mu}^{(m)})\right\|_2^2\right]
    }{
    \sum_j
    \exp\left[-\frac{1}{2\ell^2}
    \left\|S_\rho^{-1}([\hat{\mathbf{x}}_k^T,\mathbf{u}_k^T]^T-\boldsymbol{\mu}^{(j)})\right\|_2^2\right]
    },\\
    \hat{\mathbf{x}}_{k+1}
    &= \sum_{m=1}^{M} w_m(\hat{\mathbf{x}}_k)\hat{\mathbf{x}}_{k+1}^{(m)}.
    \label{eq:code_model_stitching}
\end{align}
The default implementation uses $M=6$, ridge $\lambda=10^{-7}$, and chooses
the scheduling length scale $\ell$ from the median distance between local-model
centers in normalized state-control space. Validation is an open-loop rollout
from the validation initial condition using only the pilot-command history. The
same trim-grid-trained stitched model is evaluated against each validation
dataset, including the trim-grid validation set where it is expected to perform
well and the aggressive datasets where it tests cross-regime extrapolation. This is
\texttt{run\_model\_stitching}.

\paragraph{Subspace-Hankel.}
For a past window $p$, the code builds lagged features
\[
    \mathbf{f}_k =
    [\mathbf{x}_{k-p+1}^T,\ldots,\mathbf{x}_k^T,
     \mathbf{u}_{k-p+1}^T,\ldots,\mathbf{u}_k^T]^T.
\]
After centering, an SVD basis $B$ is formed from the first $n_z$ right singular
vectors and $\mathbf{z}_k=(\mathbf{f}_k-\bar{\mathbf{f}})^TB$. The identified
latent model is
\begin{align}
    \mathbf{z}_{k+1} &\approx [1,\mathbf{z}_k^T,\mathbf{u}_k^T]A_z,\\
    \mathbf{x}_{k} &\approx [1,\mathbf{z}_k^T,\mathbf{u}_k^T]C_z.
    \label{eq:code_subspace}
\end{align}
The current defaults are $p=12$ and $n_z=6$. This is
\texttt{run\_subspace\_hankel}.

\paragraph{Koopman-EDMD.}
The EDMD row uses the same library $\Theta(\mathbf{x},\mathbf{u})$ as
\eqref{eq:code_sindy_library}, but fits a discrete lifted map rather than
continuous-time derivatives. With standardized library columns
$\bar{\Theta}$, the code solves
\begin{equation}
    \widehat{K}_x
    =
    \arg\min_K
    \|\mathbf{X}_{k+1}-\bar{\Theta}(\mathbf{x}_k,\mathbf{u}_k)K\|_F^2
    +\lambda\|K\|_F^2.
    \label{eq:code_koopman_edmd}
\end{equation}
Validation then iterates
\[
    \hat{\mathbf{x}}_{k+1}
    =
    \bar{\Theta}(\hat{\mathbf{x}}_k,\mathbf{u}_k)\widehat{K}_x.
\]
This is a compact EDMD-style lifted predictor implemented in
\texttt{fit\_koopman\_edmd} and \texttt{run\_koopman\_edmd}. It should be
interpreted as a predictive surrogate, not as an aerodynamic-parameter
estimator.

\paragraph{Frequency-Welch and Frequency-Stitching.}
The frequency-domain method fits the same continuous-time linear form as a
local state-space model,
\begin{equation}
    \dot{\mathbf{x}}=\mathbf{c}+A\mathbf{x}+B\mathbf{u}.
\end{equation}
For each Hann-windowed segment, the zero-mean Fourier coefficient of
$[\mathbf{x}^T,\mathbf{u}^T]^T$ is $\mathbf{Z}(\omega)$. The target is
$j\omega\mathbf{X}(\omega)$, and the code solves
\begin{equation}
    \widehat{G}
    =
    \arg\min_G
    \sum_{\omega\in\Omega}
    w_\omega^2
    \left\|
    j\omega\mathbf{X}(\omega)
    -
    \mathbf{Z}(\omega)^TG
    \right\|_2^2
    +\lambda\|G\|_F^2,
    \label{eq:code_frequency}
\end{equation}
then sets
$\widehat{\mathbf{c}}=\overline{\dot{\mathbf{x}}}
-\overline{[\mathbf{x}^T,\mathbf{u}^T]}\widehat{G}$.
The retained bins satisfy the coded frequency and coherence gates. This is
\texttt{fit\_frequency\_linear\_model}; it is a Welch/ETFE-inspired baseline,
not CIFER.
The Frequency-Stitching row applies the same estimator locally on the trim-grid
groups used for model stitching. For local model $i$,
\[
    \dot{\mathbf{x}}_i
    =
    \widehat{\mathbf{c}}_i
    +
    \widehat{A}_i\mathbf{x}
    +
    \widehat{B}_i\mathbf{u},
\]
and validation uses
\[
    \dot{\mathbf{x}}
    =
    \sum_i w_i(\mathbf{x},\mathbf{u})\dot{\mathbf{x}}_i,
    \qquad
    w_i =
    \frac{\exp\{-\frac{1}{2}\|S^{-1}([\mathbf{x}^T,\mathbf{u}^T]^T-\boldsymbol{\mu}_i)\|^2/\ell^2\}}
    {\sum_j \exp\{-\frac{1}{2}\|S^{-1}([\mathbf{x}^T,\mathbf{u}^T]^T-\boldsymbol{\mu}_j)\|^2/\ell^2\}}.
\]
This row therefore represents a local frequency-domain approximation, not a
new nonlinear frequency-response estimator.

\paragraph{SINDy and stepwise symbolic regression.}
Both sparse-regression methods use the same library
\begin{multline}
\Theta(\mathbf{x},\mathbf{u})=
[1,V,\alpha,\gamma,Q,T,\delta_e,V\alpha,V\gamma,\alpha Q,Q\delta_e,
T\alpha,\\
\alpha^2,\gamma^2,Q^2,\delta_e^2,\sin\alpha,\sin\gamma,\cos\alpha,\cos\gamma].
\label{eq:code_sindy_library}
\end{multline}
The implemented SINDy row does not rediscover the full vector field from
scratch. The equations of motion are inverted on the training samples to form
inferred coefficients
$\widehat{\mathbf{c}}=[\widehat{C}_L,\widehat{C}_D,\widehat{C}_m]$. Each
coefficient is then fit with a protected low-order aircraft library and a sparse
residual library:
\begin{align}
    \widehat{C}_L &\approx
        [1,\alpha]\beta_L+\Theta_{L,\mathrm{res}}\xi_L,\\
    \widehat{C}_D &\approx
        [1,\alpha,\alpha^2]\beta_D+\Theta_{D,\mathrm{res}}\xi_D,\\
    \widehat{C}_m &\approx
        [1,\alpha,Q,\delta_e]\beta_m+\Theta_{m,\mathrm{res}}\xi_m.
        \label{eq:code_sindy_structured}
\end{align}
The protected $\beta$ terms are not thresholded. The residual terms are
standardized, initialized by ridge regression, and then refit by sequential
hard-thresholded least squares for eight iterations:
\begin{equation}
    \xi_{\mathcal{A},j}
    =
    \arg\min_\beta
    \|\widehat{c}_j-\Theta_{\mathcal{P},j}\beta_{\mathcal{P}}
      -\Theta_{\mathcal{A},j}\beta\|_2^2
    +\lambda\|\beta\|_2^2,
    \qquad
    \mathcal{A}=\{i:|\xi_{ij}|\ge\tau\}.
\end{equation}
Validation integrates the flight-dynamics ODE using the learned coefficient
functions from \eqref{eq:code_sindy_structured}, initialized from the validation
initial condition and driven only by the pilot-command history. Thus SINDy is
compared as a sparse aerodynamic-closure method rather than as an unconstrained
black-box state-derivative model.

The stepwise symbolic baseline greedily adds terms that reduce
\[
    \log(\mathrm{MSE})+
    \rho\,n_{\mathrm{terms}}\log(N)/N
\]
and then refits ridge coefficients on the selected support. These correspond to
\texttt{fit\_sindy} and \texttt{fit\_symbolic\_stepwise}.

\paragraph{UDE residual.}
The UDE residual target is the measured derivative error relative to the
nominal model,
\begin{equation}
    \mathbf{r}_k
    =
    \dot{\widetilde{\mathbf{x}}}_k
    -
    f_{\mathrm{nom}}(\widetilde{\mathbf{x}}_k,\mathbf{u}_k;
    \boldsymbol{\theta}_{\mathrm{nom}}).
\end{equation}
A tanh MLP $\mathcal{N}_\phi$ is trained with normalized mean-square error, and
validation rollouts use
\begin{equation}
    \dot{\hat{\mathbf{x}}}
    =
    f_{\mathrm{nom}}(\hat{\mathbf{x}},\mathbf{u};
    \boldsymbol{\theta}_{\mathrm{nom}})
    +g_{\mathrm{UDE}}\mathcal{N}_\phi(\hat{\mathbf{x}},\mathbf{u}).
    \label{eq:code_ude}
\end{equation}
The source-aware rollout gain is $g_{\mathrm{UDE}}=1$ for direct-state runs and
$0.1$ for mocap-derived runs unless overridden. This is \texttt{run\_ude}.

\paragraph{PINN-style coefficient closure.}
The implemented PINN-style closure does not enforce CFD or an assumed complete
aerodynamic model. Instead it uses \eqref{eq:code_inferred_coefficients} to form
training targets
\begin{equation}
    \Delta\widehat{\mathbf{C}}_k
    =
    \widehat{\mathbf{C}}_k
    -
    \mathbf{C}_{\mathrm{nom}}(\widetilde{\mathbf{x}}_k,\mathbf{u}_k;
    \boldsymbol{\theta}_{\mathrm{nom}}),
\end{equation}
trains an MLP $\mathcal{N}_{C,\phi}$, and rolls out
\begin{equation}
    \dot{\hat{\mathbf{x}}}
    =
    f\!\left(
    \hat{\mathbf{x}},\mathbf{u},
    \mathbf{C}_{\mathrm{nom}}(\hat{\mathbf{x}},\mathbf{u})
    +g_{\mathrm{PINN}}\mathcal{N}_{C,\phi}(\hat{\mathbf{x}},\mathbf{u})
    \right).
    \label{eq:code_pinn_closure}
\end{equation}
The default gains match the UDE defaults. This is
\texttt{run\_pinn\_closure}.

\paragraph{GP/RBF coefficient closure.}
The GP-style row uses the same inferred coefficient-residual target
$\Delta\widehat{\mathbf{C}}_k$ as the PINN-style closure, but replaces the MLP
with a sparse radial-basis expansion. After standardizing
$\mathbf{z}_k=[\mathbf{x}_k^T,\mathbf{u}_k^T]^T$, the code selects $M$ centers
$\mathbf{c}_m$ from the training data and forms
\begin{equation}
    \phi_m(\mathbf{z})
    =
    \exp\!\left(
        -\frac{\|\bar{\mathbf{z}}-\mathbf{c}_m\|_2^2}{2\ell^2}
    \right),
    \qquad
    \Phi=[1,\phi_1,\ldots,\phi_M].
\end{equation}
The coefficient residual weights are obtained from
\begin{equation}
    \widehat{W}_{\mathrm{GP}}
    =
    \arg\min_W
    \|\Delta\widehat{\mathbf{C}}-\Phi W\|_F^2
    +\lambda\|W\|_F^2.
    \label{eq:code_gp_closure}
\end{equation}
This finite-basis kernel-ridge model is a sparse approximation to an RBF
Gaussian-process coefficient surrogate. It provides a classical
uncertainty-aware surrogate family in the benchmark taxonomy, although the
current implementation reports only mean predictions and training loss rather
than posterior variance. This is \texttt{run\_gp\_coeff\_closure}.

\paragraph{Hidden-control UDE and hidden-elevator PINN.}
For hidden-control learning, the code first infers the effective input that
would make the nominal model match the measured derivatives:
\begin{align}
    T_{\mathrm{eff},k}
    &=
    \frac{
    m\dot{\widetilde{V}}_k
    +D_{\mathrm{nom},k}
    +mg\sin\widetilde{\gamma}_k}
    {\max(\cos\widetilde{\alpha}_k,0.25)},\\
    \delta_{e,\mathrm{eff},k}
    &=
    \frac{
    J_y\dot{\widetilde{Q}}_k/(\bar{q}_kS)
    -C_{M0}-C_{M\alpha}\widetilde{\alpha}_k-C_{MQ}\widetilde{Q}_k}
    {C_{M\delta_e}}.
\end{align}
The target correction is
$\Delta\mathbf{u}_k=[T_{\mathrm{eff},k},\delta_{e,\mathrm{eff},k}]^T
-\mathbf{u}_k$, clipped to $[-0.45,0.45]$ N in throttle and
$[-0.22,0.22]$ rad in elevator. The hidden-control UDE learns both channels and
rolls out
\[
    \dot{\hat{\mathbf{x}}}
    =
    f_{\mathrm{nom}}\!\left(
    \hat{\mathbf{x}},
    \mathrm{clip}(\mathbf{u}+\mathcal{N}_{u,\phi}(\hat{\mathbf{x}},\mathbf{u}))
    \right).
\]
The hidden-elevator PINN uses the same target but constrains the throttle
correction to zero and learns only $\Delta\delta_e$. These are
\texttt{run\_ude\_hidden\_control} and \texttt{run\_pinn\_hidden\_control}.

\paragraph{Supervised coefficient surrogate.}
The neural coefficient surrogate is an oracle upper-bound row. It trains the
same MLP architecture directly on simulator labels
\[
    \Delta\mathbf{C}_{\mathrm{aero},k}
    =
    \mathbf{C}_{\mathrm{true},k}-\mathbf{C}_{\mathrm{nom},k},
\]
which are not available in real flight data. This is
\texttt{run\_supervised\_coeff\_surrogate}.

\paragraph{Implementation audit map.}
Table~\ref{tab:code_audit_map} records where each mathematical statement above
is implemented. The table is deliberately included in the paper because the
benchmark is meant to be reproducible, not merely illustrative.

\begingroup
\scriptsize
\setlength{\tabcolsep}{2pt}
\begin{longtable}{@{}p{0.18\linewidth}p{0.22\linewidth}p{0.54\linewidth}@{}}
\caption{Code-to-equation audit map for the current benchmark implementation.}
\label{tab:code_audit_map}\\
\toprule
Benchmark component & Equations in this document & Python implementation \\
\midrule
\endfirsthead
\toprule
Benchmark component & Equations in this document & Python implementation \\
\midrule
\endhead
Simulator plant and residual aerodynamics &
\eqref{eq:aero_loads}--\eqref{eq:eom} &
\texttt{methods/simulation/longitudinal.py}: \texttt{nominal\_coefficients},
\texttt{nonlinear\_residual\_coefficients}, \texttt{dynamics},
\texttt{rk4\_step}. \\
SAFE-on hidden input model &
\eqref{eq:safe_loop} &
\texttt{proprietary\_autopilot\_command}; estimator analogue in
\texttt{casadi\_controller\_command}. \\
Actuator lag/rate limits &
\eqref{eq:Tact}--\eqref{eq:eact} &
\texttt{actuator\_step}; constants checked against the current code. \\
Mocap measurement and reconstruction &
\eqref{eq:code_derivative_pipeline} &
\texttt{mocap\_from\_state}, \texttt{derive\_state\_from\_mocap},
\texttt{smooth\_trials}. \\
Equation-error LS &
\eqref{eq:code_inferred_coefficients}--\eqref{eq:code_equation_error} &
\texttt{infer\_coefficients}, \texttt{fit\_equation\_error}. \\
OEM-SS and mocap-output OEM &
\eqref{eq:code_oem_ss}--\eqref{eq:code_oem_mocap} &
\texttt{run\_oem\_casadi}, \texttt{run\_oem\_mocap\_output\_casadi}. \\
Variational-Mocap &
\eqref{eq:code_variational_mocap} &
\texttt{run\_variational\_mocap\_output\_casadi}. \\
OEM-HiddenController &
\eqref{eq:code_oem_hidden} &
\texttt{run\_oem\_hidden\_controller}. \\
EKF-ParamID &
\eqref{eq:code_ekf} &
\texttt{make\_casadi\_rk4\_step\_parameter\_jacobian}, \texttt{run\_filter\_error\_ekf}. \\
Linear, stitched, and subspace models &
\eqref{eq:code_linear_ss}--\eqref{eq:code_subspace} &
\texttt{run\_linear\_state\_space}, \texttt{run\_model\_stitching},
\texttt{run\_subspace\_hankel}. \\
Koopman-EDMD &
\eqref{eq:code_koopman_edmd} &
\texttt{fit\_koopman\_edmd}, \texttt{run\_koopman\_edmd}. \\
Frequency-Welch and Frequency-Stitching &
\eqref{eq:code_frequency} &
\texttt{fit\_frequency\_linear\_model}, \texttt{run\_frequency\_linear},
\texttt{run\_frequency\_stitching}. \\
SINDy and symbolic stepwise &
\eqref{eq:code_sindy_library} &
\texttt{sindy\_library}, \texttt{fit\_sindy}, \texttt{fit\_symbolic\_stepwise}. \\
UDE, PINN closure, hidden-control neural rows &
\eqref{eq:code_ude}--\eqref{eq:code_gp_closure} &
\texttt{run\_ude}, \texttt{run\_pinn\_closure},
\texttt{run\_gp\_coeff\_closure}, \texttt{infer\_input\_correction}, \texttt{run\_ude\_hidden\_control},
\texttt{run\_pinn\_hidden\_control}. \\
\bottomrule
\end{longtable}
\endgroup

\subsection*{Shared Benchmark Results}
The benchmark is generated by a reproducible Python workflow that writes the
method-comparison tables, diagnostic figures, and LaTeX-ready artifacts directly
into the \texttt{latex} directory. The main text emphasizes graphics because the
raw tables are too verbose to be the primary interpretation tool. The full
sorted tables are retained in Appendix~\ref{app:benchmark_tables} for audit and
reproduction. The validation score is a normalized aggregate trajectory error,
so lower values indicate better agreement with the simulated truth. Training
time and rollout time are reported separately because the computational burden
differs across method families: CasADi/IPOPT methods spend most of their time in
nonlinear programming, neural methods train quickly on the GPU but require
repeated rollout evaluations, and linear or regression-based methods fit almost
immediately but can be sensitive to measurement preprocessing.

\begin{figure}[!htbp]
    \centering
    \includesvg[width=0.98\linewidth]{generated_benchmark_problem_matrix}
    \caption{Benchmark interpretation axes. The paper evaluates both direct-state and motion-capture-style measurement assumptions, and it separately studies SAFE-off and SAFE-on hidden-input dynamics.}
    \label{fig:benchmark_problem_matrix}
\end{figure}

Figure~\ref{fig:benchmark_maneuver_overview} shows representative validation
maneuvers for the benchmark datasets. For each dataset, the plotted trajectory
is the validation trial with the largest absolute angle of attack, so the figure
shows the intended excitation envelope rather than an arbitrary mild trial. The
aggressive sine-sweep and aggressive nonlinear cases deliberately add large
elevator pulses, throttle excursions, and high-$\alpha$ dwell segments to leave
the near-trim region. The SAFE-on variants receive similarly aggressive pilot
commands, but the hidden pitch/rate recovery logic reshapes the actuator input
and suppresses much of the actual high-$\alpha$ response.

\begin{figure}[!htbp]
    \centering
    \includesvg[width=0.98\linewidth]{generated_benchmark_maneuver_overview}
    \caption{Representative validation maneuvers. The path panel shows inertial motion-capture position $p_x$ versus $p_z$ after subtracting the initial position. The remaining panels show airspeed, angle of attack, pitch attitude, thrust command/realized input, and elevator command/realized input for the same trials. Dotted input traces are external pilot commands; solid or dashed traces are the actuator inputs that drive the plant.}
    \label{fig:benchmark_maneuver_overview}
\end{figure}

The current benchmark does not include full loops. The 3-DoF longitudinal model
uses $V,\alpha,\gamma,q$ with a conventional small-disturbance-style flight-path
angle envelope, and the simulator rejects trajectories outside
$4\leq V\leq34$ m/s, $|\gamma|\leq1.20$ rad, $|\alpha|\leq0.75$ rad, and
$|q|\leq3.50$ rad/s. These limits are deliberately wider than the nominal pylon
racing envelope so that aggressive validation trajectories can include
post-stall pull-ups, rapid pitch excursions, and steep climbs or descents. A
true loop would require $\gamma$ and pitch attitude to pass through a full
rotation, with wrapped attitude coordinates and an aerodynamic model valid for
inverted, post-stall, and possibly negative-lift flight. That would be a useful
future stress case, but it should be treated as a separate aerobatic benchmark
rather than mixed into the present indoor trainer identification study.

\begin{figure}[!htbp]
    \centering
    \includesvg[width=0.98\linewidth]{generated_method_score_heatmap_direct}
    \caption{Direct-state open-loop validation trajectory error by method and benchmark condition. Lower values are better. All ranked methods use validation measurements only for initialization and scoring, not for feedback during validation rollout. The leftmost column gives the per-method mean score across the direct-state benchmark conditions.}
    \label{fig:method_score_heatmap_direct}
\end{figure}

\begin{figure}[!htbp]
    \centering
    \includesvg[width=0.98\linewidth]{generated_method_score_heatmap_mocap}
    \caption{Motion-capture-derived open-loop validation trajectory error by method and benchmark condition. Lower values are better. The leftmost column gives the per-method mean score across the mocap-derived benchmark conditions, separating the measurement-reconstruction problem from the direct-state comparison in Fig.~\ref{fig:method_score_heatmap_direct}.}
    \label{fig:method_score_heatmap_mocap}
\end{figure}

For the direct-state case, the supervised coefficient surrogate gives the best
score because it uses synthetic labels for the hidden aerodynamic residuals; it
is therefore best interpreted as an upper-bound surrogate rather than a
flight-data-only method. The UDE, PINN-style coefficient-closure, and GP/RBF
coefficient-closure models test the flight-data-only version of the same idea:
retain the rigid-body equations while learning missing force and moment closures
from residuals. This ranking now changes as intended with maneuver severity.
In the near-trim open-loop direct-state case, many predictors work well and
Linear-SS remains a reasonable local model ($J_{\mathrm{val}}\approx0.041$).
In the aggressive direct-state case, the learned aerodynamic-residual family
separates from the global linear model: the supervised coefficient surrogate,
UDE residual, and PINN-style coefficient closure obtain
$J_{\mathrm{val}}\approx0.022$, $0.026$, and $0.028$, respectively, while
Linear-SS is about $0.19$. Thus the aggressive benchmark no longer supports the
interpretation that a single global affine state-space model is sufficient once
the validation set spans stall onset, recovery, and speed variation. The
Koopman-EDMD row tests a different surrogate role, namely a lifted one-step
predictor with no aerodynamic parameter interpretation. The EKF-ParamID row
uses filtering only during training; its validation score is a frozen-parameter
open-loop rollout driven by the pilot command sequence.

The mocap-derived comparison shows a different ranking. Methods that depend
strongly on differentiated states, especially equation-error, SINDy, GP/RBF
coefficient closure, PINN-style closure, UDE residual learning, and the compact
Frequency-Welch baseline, degrade when reconstructed $V,\alpha,\gamma,q$
inherit numerical differentiation noise and smoothing lag from position and
attitude. In the aggressive mocap-derived case, EKF-ParamID, mocap-output OEM,
OEM-SS, and Linear-SS cluster in the range
$J_{\mathrm{val}}\approx0.13$--$0.20$, while several nonlinear residual methods
that were strongest with direct states are no longer dominant. This should not
be read as evidence that the aircraft is globally linear; rather, it indicates
that the preprocessing pipeline has become the bottleneck and can hide the value
of nonlinear aerodynamic closures. Open-loop grey-box and learned-residual
models expose the cost of using reconstructed states as if they were direct
aerodynamic measurements. This is precisely the practical issue motivating
mocap-output OEM and latent-state filter or variational formulations.

Tables~\ref{tab:open_loop_shared_method_comparison_direct} and
\ref{tab:open_loop_safe_shared_method_comparison_direct} isolate the SAFE-off
versus SAFE-on comparison for the same randomized open-loop pilot-command
excitation. In the SAFE-off case, the main unobserved input effect is actuator
lag between $\mathbf{u}_{\mathrm{cmd}}$ and $\mathbf{u}_{\mathrm{act}}$. In the
SAFE-on case, the same observed pilot command is additionally passed through the
hidden pitch-attitude/rate loop in \eqref{eq:safe_loop}. This comparison is the
closest numerical analog to flying the off-the-shelf aircraft with the
proprietary stabilizer disabled or enabled.

The aggressive nonlinear maneuver is a separate stress case rather than a
replacement for the local open-loop maneuver. It adds faster elevator doublets,
larger pull-up and pushover pulses, and dedicated high-$\alpha$ stall probes
with long elevator holds. Its purpose is to leave the near-trim region where the
coded nonlinear aerodynamic residuals can be approximated by a local affine
model and to force the SAFE-off validation set through the lift-rolloff and
drag-rise portion of the hidden RC-aircraft polar. The resulting ranking is
therefore a test of nonlinear aerodynamic recovery and extrapolation, not only
of short-horizon in-distribution trajectory prediction.

The SAFE-on aggressive and recovery-probe cases add a second effect: the
external pilot command is no longer the input that directly drives the plant.
For direct-state validation, PINN-style coefficient closure, UDE residual
learning, Koopman-EDMD, and model stitching all improve substantially over the
nominal model, while Linear-SS remains a useful but not leading local predictor.
For mocap-derived SAFE-on validation, the hidden-controller OEM and OEM-SS rows
become more competitive because they either include an oracle hidden-controller
structure or optimize latent trajectories against the measured outputs. This is
consistent with the practical interpretation of SAFE-on data: without a model of
the hidden elevator reshaping, an estimator can easily attribute controller
action to aerodynamic residuals.

\begin{figure}[!htbp]
    \centering
    \includesvg[width=0.98\linewidth,height=0.74\textheight,keepaspectratio]{generated_dataset_direct_validation_score_panels}
    \caption{Direct-state open-loop validation score panels for the archived benchmark datasets. The SAFE-off and SAFE-on panels isolate how hidden command reshaping changes the ranking even before mocap reconstruction noise is introduced. Lower score is better.}
    \label{fig:dataset_direct_validation_score_panels}
\end{figure}
\FloatBarrier

Tables~\ref{tab:safe_loop_shared_method_comparison_direct} and
\ref{tab:safe_loop_shared_method_comparison_mocap} show the SAFE-loop recovery
probe comparison. This dataset is the most relevant to the off-the-shelf vehicle
question because the estimator sees external pilot commands while the plant is
driven by internally reshaped elevator commands. The hidden-controller OEM is
included in this table, but it is intentionally marked as an oracle-structure
case: it assumes that the analyst has guessed a parameterized SAFE-like pitch
law. That assumption is optimistic. In real proprietary firmware, saturations,
rate limits, and panic/recovery gates are difficult to identify unless the test
data visit those regions, and their absence from ordinary lap data can make the
controller appear as ordinary aerodynamic or actuator model error.

Figure~\ref{fig:shared_train_time_accuracy_tradeoff} places the same results on a cost-error plot. The horizontal axis reports training or nonlinear-solve wall time, while the vertical axis reports $J_{\mathrm{val}}$, the mean normalized state RMSE defined above; lower vertical position is better. The detailed appendix tables also report process CPU time, CUDA event training time, and peak CUDA memory for PyTorch rows when a GPU is used. This distinction matters when runs are parallelized: wall time measures user wait time, CPU time measures host work consumed by each process, and GPU event time isolates the neural training kernels. In the direct-state case, neural residual and coefficient-closure approaches achieve low error with short GPU training time, whereas CasADi/IPOPT OEM methods spend more time in nonlinear optimization. In the mocap-derived case, the dominant issue is no longer training cost but measurement preprocessing: several fast derivative-based methods become inaccurate because the reconstructed aerodynamic states are poor dependent variables.

\begin{figure}[H]
    \centering
    \includesvg[width=0.98\linewidth]{generated_shared_train_time_accuracy_tradeoff}
    \caption{Training-time versus validation-error tradeoff for the near-trim open-loop benchmark. The vertical axis is $J_{\mathrm{val}}$, the mean of the range-normalized RMSEs for $V,\alpha,\gamma$, and $q$ over all validation samples. Methods with $J_{\mathrm{val}}>1.0$ are treated as failed rollouts and summarized in the callout so the useful region remains visible. Marker area increases with validation rollout time. The red dashed horizontal line gives the nominal-model baseline, which has no fitted training solve. Lower vertical position means lower trajectory error and therefore better validation performance.}
    \label{fig:shared_train_time_accuracy_tradeoff}
\end{figure}

\IfFileExists{fig/generated_aggressive_shared_train_time_accuracy_tradeoff.svg}{%
\begin{figure}[H]
    \centering
    \includesvg[width=0.98\linewidth]{generated_aggressive_shared_train_time_accuracy_tradeoff}
    \caption{Training-time versus validation-error tradeoff for the aggressive nonlinear benchmark. Methods with $J_{\mathrm{val}}>1.0$ are summarized as failed rollouts in the callout. This case uses stratified initial conditions and larger elevator/throttle excitation to drive stall onset, drag rise, speed variation, pull-up, and recovery regimes. It should be interpreted as the primary nonlinear aerodynamic stress test, while Fig.~\ref{fig:shared_train_time_accuracy_tradeoff} is the local near-trim comparison.}
    \label{fig:aggressive_train_time_accuracy_tradeoff}
\end{figure}
}{}

\IfFileExists{fig/generated_aggressive_safe_shared_train_time_accuracy_tradeoff.svg}{%
\begin{figure}[H]
    \centering
    \includesvg[width=0.98\linewidth]{generated_aggressive_safe_shared_train_time_accuracy_tradeoff}
    \caption{Training-time versus validation-error tradeoff for the aggressive benchmark with the hidden SAFE-like controller enabled. Methods with $J_{\mathrm{val}}>1.0$ are summarized as failed rollouts in the callout. The estimator receives the same external pilot-command channel, while the simulator is driven by an internally reshaped elevator command. This figure isolates the combined effect of nonlinear aerodynamics and unknown closed-loop input dynamics.}
    \label{fig:aggressive_safe_train_time_accuracy_tradeoff}
\end{figure}
}{}

Figure~\ref{fig:shared_validation_score_comparison} shows the same validation scores as a ranked log-scale score plot split by measurement assumption.

\begin{figure}[!htbp]
    \centering
    \includesvg[width=0.98\linewidth]{generated_shared_validation_score_comparison}
    \caption{Open-loop validation trajectory score for the shared 3-DoF longitudinal benchmark. The score is the mean state NRMSE $J_{\mathrm{val}}$, so lower is better. Direct-state rows use simulated noisy $V,\alpha,\gamma,q$ measurements; mocap rows use states reconstructed from 100 Hz position and attitude measurements.}
    \label{fig:shared_validation_score_comparison}
\end{figure}

The EKF-ParamID row should be interpreted as a training-time parameter
estimator, not as a validation-time state estimator. The filter uses training
measurements to update the augmented state and parameter covariance, then
freezes the learned parameters for validation. This avoids the unfair
comparison that would occur if validation measurements were assimilated at
every time step. If a separate state-estimation diagnostic is desired, it
should be reported separately from the open-loop model-identification rankings.

\begin{figure}[!htbp]
    \centering
    \includesvg[width=0.98\linewidth]{generated_shared_validation_trajectory_overlay}
    \caption{Representative validation trajectory overlay for the best-scoring methods from the shared benchmark. The plot highlights the difference between direct-state tracking and methods that are sensitive to derivative reconstruction and model-form error.}
    \label{fig:shared_validation_trajectory_overlay}
\end{figure}

The observation-rate study in Fig.~\ref{fig:observation_rate_state_reconstruction} isolates the effect of lowering the motion-capture observation rate; the full numerical table is in Appendix~\ref{app:benchmark_tables}. The degradation in reconstructed state quality is modest when going from 100 Hz to 10 Hz for slow longitudinal content, but becomes much more severe at 2 Hz, where the Nyquist frequency is only 1 Hz. At that point the short-period and actuator-driven content needed for identification can be attenuated or aliased before any estimator sees the data.

\begin{figure}[!htbp]
    \centering
    \includesvg[width=0.82\linewidth]{generated_observation_rate_state_reconstruction}
    \caption{State-reconstruction error obtained by decimating the 100 Hz mocap stream to lower observation rates before smoothing and finite differencing.}
    \label{fig:observation_rate_state_reconstruction}
\end{figure}

The Fisher-information diagnostics in Appendix~\ref{app:benchmark_tables} provide a local uncertainty view for the equation-error aerodynamic parameter fit. These values should not be interpreted as full Bayesian posterior intervals; they are local CRLB-style diagnostics around a fitted parameter vector under the assumed noise model. High parameter correlations indicate directions in the aerodynamic parameter space that are weakly separated by the available maneuver data, reinforcing the need for excitation design and uncertainty-aware validation.

\begin{figure}[!htbp]
    \centering
    \includesvg[width=0.95\linewidth]{generated_shared_pinn_coeff_residual_validation}
    \caption{PINN-style coefficient-closure residual comparison on a validation trajectory. This diagnostic tests whether the learned closure recovers the hidden nonlinear aerodynamic coefficient increments rather than merely matching states.}
    \label{fig:shared_pinn_coeff_residual_validation}
\end{figure}

\begin{figure}[!htbp]
    \centering
    \includesvg[width=0.95\linewidth]{generated_shared_frequency_validation_diagnostic}
    \caption{Frequency-domain diagnostic from the validation maneuver. The coherence and frequency-response estimates are useful for checking excitation bandwidth, but ordinary maneuver data are not a substitute for a dedicated multisine or sine-sweep frequency-identification experiment.}
    \label{fig:shared_frequency_validation_diagnostic}
\end{figure}

\FloatBarrier

\section{Supplemental Method Formulations}
This section keeps equation-level details for implemented method families whose
benchmark results are generated by the repository. Legacy standalone figures
and tables that are not regenerated by the current pipeline have been removed.

\subsection{Output Error Method}
This formulation belongs to the class of prediction error methods (PEM), where the objective is to minimize the discrepancy between model-predicted outputs and measured data. The Output Error (OE) method is a special case that assumes deterministic system dynamics, with all uncertainty entering through measurement noise.

Starting from the nonlinear system,
\begin{align}
\dot{x}(t) &= f\big(x(t),u(t);\theta_p\big),\\
y(t) &= h\big(x(t);\theta_p\big) + v(t),
\end{align}
the predicted output $\hat{y}(t,\theta_p)$ is obtained by numerically integrating the state equations for a given parameter vector $\theta_p$ (see Section~\ref{sec:prob}).

Assuming additive Gaussian measurement noise,
\begin{align}
y_k = \hat{y}_k(\theta_p) + \varepsilon_k, \qquad \varepsilon_k \sim \mathcal{N}(0, R),
\end{align}
the maximum likelihood estimation (MLE) problem becomes
\begin{align}
J_{\text{MLE}}(\theta_p, R) =
\frac{1}{2} \sum_{k=1}^{N} 
\left(y_k - \hat{y}_k(\theta_p)\right)^\top R^{-1} 
\left(y_k - \hat{y}_k(\theta_p)\right)
+ \frac{N}{2} \log \det R.
\end{align}
In this work, $R$ is assumed diagonal and estimated jointly with $\theta_p$ via a log-parameterization, resulting in adaptive weighting across output channels.

The resulting nonlinear program is constructed symbolically using CasADi and solved using IPOPT, with state trajectories generated through a fixed-step fourth-order Runge--Kutta (RK4) integrator embedded within the optimization.

Two OE formulations are considered: single shooting (SS) and multiple shooting (MS).

\subsubsection{Single--Shooting Formulation}

In the single--shooting implementation, the full state trajectory is generated by
integrating the nominal longitudinal model over the complete experiment horizon.
The aerodynamic parameter vector, the initial state, and the measurement-noise
standard deviations are estimated jointly. The decision vector is
\[
    w_{\mathrm{SS}}
    =
    \begin{bmatrix}
        \boldsymbol{\theta}_p^\top &
        \mathbf{x}_0^\top &
        \log\boldsymbol{\sigma}^\top
    \end{bmatrix}^{\!\top}.
\]
Starting from the optimized initial condition $\mathbf{x}_0$, the trajectory is
propagated using a fixed-step fourth-order Runge--Kutta map,
\[
    \hat{\mathbf{x}}_{k+1}
    =
    \Phi(\hat{\mathbf{x}}_k,\mathbf{u}_k;\boldsymbol{\theta}_p,\Delta t),
    \qquad k=0,\ldots,N-1.
\]
In the oracle state-measurement comparison, the predicted output is
taken as $\hat{\mathbf{y}}_k=\hat{\mathbf{x}}_k$. In the motion-capture comparison,
the same latent trajectory is instead passed through the measurement model
\[
\hat{\mathbf{y}}_{\mathrm{mocap},k}
=h(\hat{\mathbf{x}}_k)
=
\begin{bmatrix}
\hat{p}_{x,k} & \hat{p}_{z,k} & \hat{\alpha}_k+\hat{\gamma}_k
\end{bmatrix}^{T},
\]
where $\hat{p}_x$ and $\hat{p}_z$ are obtained by integrating
$\dot{p}_x=\hat{V}\cos\hat{\gamma}$ and
$\dot{p}_z=\hat{V}\sin\hat{\gamma}$. This output-error formulation avoids
using numerically differentiated mocap data as the dependent variable. The negative log-likelihood
objective, assuming independent Gaussian noise in each channel, is
\[
    J_{\mathrm{SS}}
    =
    \sum_{j=1}^{4} N\log\sigma_j
    +
    \frac{1}{2}
    \sum_{k=0}^{N-1}
    \sum_{j=1}^{4}
    \frac{\left(\hat{y}_{k,j}-y_{k,j}\right)^2}{\sigma_j^2}.
\]
Thus, the single--shooting OEM problem is
\[
    \min_{\boldsymbol{\theta}_p,\,\mathbf{x}_0,\,\log\boldsymbol{\sigma}}
    J_{\mathrm{SS}}.
\]
While this formulation is compact, the entire predicted trajectory relies on a single integration from the initial state, rendering the optimization sensitive to parameter initialization and prone to error accumulation over long horizons. Consequently, the Equation Error Method is often employed as a precursor to provide a robust initial guess for the OEM. The single-shooting output error approach remains the current industry standard, implemented in widely used toolkits such as SIDPAC \cite{Morelli2002SIDPAC}.

\subsection{Variational System Identification}
Variational system identification occupies the space between deterministic OEM and filter-error methods. For a stochastic continuous-time model,
\[
d\mathbf{x}(t)=f(\mathbf{x}(t),\mathbf{u}(t);\theta_p)\,dt+G(\theta_p)\,dW(t),
\qquad
\mathbf{y}_k=h(\mathbf{x}_k,\mathbf{u}_k;\theta_p)+\mathbf{v}_k,
\]
the exact likelihood requires marginalizing over the latent state path $\mathbf{x}_{0:N}$. FEM approximates this marginalization using a Kalman-filter prediction model. Variational system identification instead introduces an assumed smoothing density $q(\mathbf{x}_{0:N})$ and maximizes an evidence lower bound,
\[
\mathcal{L}(\theta_p,q)
=
\mathbb{E}_{q}
\left[
\log p_{\theta_p}(\mathbf{x}_{0:N},\mathbf{y}_{0:N})
-\log q(\mathbf{x}_{0:N})
\right],
\]
over both the model parameters and the smoothing distribution. In Dutra's aircraft formulation \citep{dutra_variational_2025}, the assumed density is chosen to preserve a Markov structure so that the objective remains sparse and tractable. This makes VI a natural candidate for the mocap-output problem considered here: the measured outputs are $p_x,p_z,\theta$, the flight states are latent, and turbulence or unmodeled disturbances can be represented as process noise rather than being forced into measurement residuals.

The role of VI in this repository benchmark is therefore different from SINDy, PINNs, or UDEs. VI improves inference for a prescribed grey-box model in the presence of process and measurement noise; it does not automatically supply a new aerodynamic closure. If the nominal coefficient structure is incomplete, VI can produce better state smoothing and uncertainty estimates, but the missing aerodynamic physics must still be represented through additional parameters, basis functions, Gaussian-process terms, or neural residuals. A natural extension of the present benchmark is therefore a variational UDE or variational coefficient-closure model, where VI handles latent-state and process-noise uncertainty while the residual model handles structural aerodynamic error.

\subsubsection{Multiple--Shooting Formulation}

In the multiple--shooting implementation \cite{bock1984multiple,jategaonkar2006flight}, the time horizon is divided into $M$
shorter windows. The aerodynamic parameters and noise standard deviations remain
global, while each window has its own optimized shooting-node state. The decision
vector is
\[
    w_{\mathrm{MS}}
    =
    \begin{bmatrix}
        \boldsymbol{\theta}_p^\top &
        \mathbf{X}_0^\top &
        \mathbf{X}_1^\top &
        \cdots &
        \mathbf{X}_{M-1}^\top &
        \log\boldsymbol{\sigma}^\top
    \end{bmatrix}^{\!\top}.
\]
Within the $i$th window, the local trajectory is initialized from $\mathbf{X}_i$
and propagated as
\[
    \hat{\mathbf{x}}_{i,k+1}
    =
    \Phi(\hat{\mathbf{x}}_{i,k},\mathbf{u}_{i,k};
    \boldsymbol{\theta}_p,\Delta t),
    \qquad
    \hat{\mathbf{x}}_{i,0}=\mathbf{X}_i.
\]
Continuity between neighboring windows is enforced through equality constraints,
\[
    \mathbf{g}_i
    =
    \hat{\mathbf{x}}_{i,\mathrm{end}}
    -
    \mathbf{X}_{i+1}
    =
    \mathbf{0},
    \qquad i=0,\ldots,M-2.
\]
The corresponding MLE problem is
\[
\begin{aligned}
    \min_{\boldsymbol{\theta}_p,\,\{\mathbf{X}_i\},\,\log\boldsymbol{\sigma}}
    \quad
    J_{\mathrm{MS}}
    &=
    \sum_{j=1}^{4} N\log\sigma_j
    +
    \frac{1}{2}
    \sum_{i=0}^{M-1}
    \sum_{k=0}^{N_i-1}
    \sum_{j=1}^{4}
    \frac{
    \left(\hat{y}_{i,k,j}-y_{i,k,j}\right)^2
    }{\sigma_j^2} \\
    \text{s.t.}\quad
    \mathbf{g}_i &= \mathbf{0},
    \qquad i=0,\ldots,M-2.
\end{aligned}
\]
This introduces additional decision variables but allows the optimizer to leverage intermediate state information, which significantly improves robustness and convergence behavior. In the current benchmark, the generated comparison figures and tables are the authoritative output-error results; no separate hand-maintained SS/MS demonstration figure is included.
\subsection{SINDy and its Variations}
Sparse Identification of Nonlinear Dynamics (SINDy) searches for governing
equations by regressing measured state derivatives on a user-defined library of
candidate nonlinear functions. For the longitudinal benchmark, the dependent
variable for a fully black-box SINDy model would be the smoothed derivative matrix
\[
    \dot{X} =
    \begin{bmatrix}
    \dot{V} & \dot{\alpha} & \dot{\gamma} & \dot{Q}
    \end{bmatrix},
\]
and the candidate library is evaluated on
$[V,\alpha,\gamma,Q,T,\delta_e]$. The standard SINDy model is
\begin{equation}
    \dot{X} = \Theta(X,U)\,\Xi,
    \label{eq:sindy_standard}
\end{equation}
where each column of $\Xi$ contains the active terms for one state equation. A
LASSO-style sparse formulation can be written as
\begin{equation}
    \Xi^\star
    =
    \operatorname*{arg\,min}_{\Xi'}
    \left\|
    \Theta(X,U)\Xi' - \dot{X}
    \right\|_F^2
    +\lambda\|\Xi'\|_1 .
    \label{eq:sindy_lasso}
\end{equation}
The implementation used for the benchmark follows the common sequential
thresholded least-squares philosophy rather than a full convex LASSO solve, but
applies it to aerodynamic coefficients rather than directly to
\eqref{eq:sindy_standard}. The low-order aircraft terms are protected in the
regression, while additional nonlinear terms are sparse. This lets SINDy learn
the nominal-like coefficient parameters and the missing nonlinear aerodynamic
terms in a single fit. The same library is also used by a forward stepwise
symbolic-regression baseline, which remains a black-box derivative fit.

For aircraft data, the practical details are as important as the sparse
optimizer. Derivatives are obtained from smoothed direct-state or
mocap-derived-state histories; the latter case is intentionally difficult
because motion-capture position and attitude must be differentiated to recover
airspeed, flight-path angle, angle of attack, and pitch rate. Hidden actuator
or SAFE-loop dynamics are also challenging because the candidate library sees
the external pilot command rather than the internal control signal. This is why
the benchmark uses structured coefficient SINDy: the known rigid-body balance is
built into the coefficient-inference step, low-order aircraft terms are retained
as protected regressors, and sparse discovery is reserved for missing nonlinear
aerodynamic terms.

\subsection{Physics-Informed and Neural-Residual Methods}
For the longitudinal benchmark considered here, a direct inverse PINN that penalizes deviations from the same fixed aerodynamic coefficient model used by OEM is not the most appropriate use of physics information. The purpose of the benchmark is to evaluate identification methods when unmodeled aerodynamic forces, moments, actuator effects, or disturbances are present. If the PINN physics residual forces the trajectory to satisfy an incomplete nominal force and moment model, the residual acts as a regularizer toward the wrong dynamics and can suppress the missing effects the method is supposed to discover.

The useful distinction is between exact structure and aerodynamic closure assumptions. Kinematic identities and rigid-body balance laws should be enforced strongly or built into the parameterization. By contrast, aerodynamic force and moment functions should remain learnable because these are precisely the terms that may change outside the nominal model. A tractable PINN formulation for this benchmark therefore uses a neural trajectory $\widehat{\mathbf{x}}_\theta(t)$ and learned aerodynamic coefficient functions,
\[
\widehat{C}_{D,\phi}=\mathcal{N}_{D,\phi}(\widehat{\alpha},\widehat{q},\delta_e,\widehat{V}),\quad
\widehat{C}_{L,\phi}=\mathcal{N}_{L,\phi}(\widehat{\alpha},\widehat{q},\delta_e,\widehat{V}),\quad
\widehat{C}_{M,\phi}=\mathcal{N}_{M,\phi}(\widehat{\alpha},\widehat{q},\delta_e,\widehat{V}),
\]
rather than treating a low-order coefficient expansion as the true physics. The corresponding dimensional forces and moments are
\[
\widehat{D}=\bar{q}S\widehat{C}_{D,\phi},\qquad
\widehat{L}=\bar{q}S\widehat{C}_{L,\phi},\qquad
\widehat{M}=\bar{q}S\widehat{C}_{M,\phi},
\]
where $\bar{q}=\frac{1}{2}\rho \widehat{V}^2$. The differential-equation residual then enforces the longitudinal rigid-body equations,
\[
\begin{aligned}
r_V &= \dot{\widehat{V}}
 - \frac{1}{m}\left(-\widehat{D}+T\cos\widehat{\alpha}-mg\sin\widehat{\gamma}\right),\\
r_\gamma &= \dot{\widehat{\gamma}}
 - \frac{1}{m\widehat{V}}\left(\widehat{L}+T\sin\widehat{\alpha}-mg\cos\widehat{\gamma}\right),\\
r_q &= \dot{\widehat{q}}-\frac{\widehat{M}}{J_y}.
\end{aligned}
\]
The pitch kinematic relationship, such as $\dot{\alpha}=q-\dot{\gamma}$ for the standard longitudinal coordinates, is not an aerodynamic discovery target and should preferably be enforced by the state parameterization itself. If it is not hard-coded, it can be included as a high-weight residual, but this is a numerical convenience rather than a claim that the kinematics are uncertain physics.

The resulting training objective is
\[
\mathcal{L}(\theta,\phi)=
\frac{\lambda_{\mathrm{data}}}{N_y}\sum_i
\left\|\mathbf{H}\widehat{\mathbf{x}}_\theta(t_i)-\mathbf{y}_i\right\|_2^2
+\frac{\lambda_{\mathrm{dyn}}}{N_f}\sum_j
\left\|\mathbf{r}_{\mathrm{dyn}}(t_j)\right\|_2^2
+\lambda_{\mathrm{reg}}\mathcal{R}_{\mathrm{aero}}(\phi),
\]
where $\mathbf{r}_{\mathrm{dyn}}=[r_V,r_\gamma,r_q]^T$. The regularization term can encode weak aerodynamic priors such as smoothness, dynamic-pressure scaling, drag positivity, boundedness, and local derivative signs near trim. These priors are much weaker than enforcing a complete aerodynamic model and are therefore more consistent with the goal of identifying unmodeled force and moment behavior.

This setup places the present benchmark in the same family as the Non-Determinant and Modified Non-Determinant PINNs of Michek et al.~\citep{michek2026flight}, but at a deliberately smaller scale. Their work learns aerodynamic coefficients for a nonlinear 6-DOF simulation using many randomized flight intervals. Here, the same conceptual role is isolated in a 3-DOF longitudinal problem with known nominal coefficients and hidden nonlinear residual terms, making it possible to compare PINN-style closure learning directly against OEM, SINDy, and UDE residual models.

A UDE formulation uses the same logic but augments a nominal coefficient model with learned residuals,
\[
\begin{aligned}
C_D &= C_{D,\mathrm{nom}}(x,u;\theta_p)+\Delta C_{D,\phi}(x,u),\\
C_L &= C_{L,\mathrm{nom}}(x,u;\theta_p)+\Delta C_{L,\phi}(x,u),\\
C_M &= C_{M,\mathrm{nom}}(x,u;\theta_p)+\Delta C_{M,\phi}(x,u).
\end{aligned}
\]
This is especially plausible for the present scope because the known model is used where it is trusted, while the neural residual captures model-form error. The identified parameters $\theta_p$ remain interpretable, and the residual terms provide a direct diagnostic for where the nominal aerodynamic model is inadequate.

This formulation is intentionally not a full aerodynamic PINN in the Navier-Stokes sense. Without geometry, flow-field states, boundary conditions, and a high-fidelity solver or data set, enforcing the Euler, RANS, or Navier-Stokes equations is not tractable in this repository benchmark. The PINN comparison is therefore limited to flight-dynamics-level grey-box learning, while high-fidelity CFD-constrained surrogate modeling is treated as future work.




\section{Method Selection Guidance}
The benchmark is designed to support method selection rather than identify a
single universally best algorithm. When direct measurements of
$V,\alpha,\gamma,Q$ are available and the primary mismatch is unknown
aerodynamics, coefficient-closure and residual-learning methods are attractive
because they preserve the rigid-body equations while learning the missing force
and moment terms. The supervised coefficient surrogate is the best possible
version of this idea in the synthetic benchmark, but it uses labels that real
flight tests do not provide; it should therefore be interpreted as an upper
bound. Flight-data-only UDE and PINN-style closure methods are more realistic
choices in this regime.

When only motion capture is available, the main difficulty changes. Algorithms
that require accurate state derivatives suffer from finite-difference noise,
smoothing lag, and loss of bandwidth. For this reason, mocap-output OEM and
variational or smoothing formulations are better
matched to the open-loop model-identification problem than methods that first reconstruct
$V,\alpha,\gamma,Q$ and then differentiate them. The direct-state benchmark is
still useful, but it should be read as an algorithm diagnostic, not as the
expected performance for the PURT pylon-racing measurement setup.

When SAFE is enabled, the estimator observes the external pilot command while
the plant responds to an internally reshaped elevator command. In this setting,
a purely aerodynamic residual model can fit trajectories while misattributing
hidden control action to aerodynamic effects. A hidden-input UDE or a
parameterized hidden-controller OEM is more conceptually appropriate. The OEM
hidden-controller row is deliberately optimistic because it assumes the analyst
has guessed the controller structure. Real proprietary controllers may contain
mode switches, saturations, rate limits, and recovery gates that are only
identifiable if the flight test excites those regions. Therefore, ordinary lap
data may be sufficient to learn a local predictive model but insufficient to
reverse-engineer the controller.

For frequency-domain identification, dedicated excitation matters. The compact
Frequency-Welch baseline is useful for checking bandwidth and coherence, but it
is not a substitute for a purpose-built CIFER-style campaign with carefully
designed multisines, sweeps, input separation, and coherence validation. The
sine-sweep datasets are included to make this distinction explicit.

The practical recommendation is therefore problem dependent. Use mocap-output
OEM or variational/filter-error methods when the measurement system is the
dominant limitation. Use UDE or PINN-style coefficient closures when the
measurement states are reliable but the aerodynamic model is incomplete. Use
GP/RBF coefficient closures when a kernel surrogate with smoother interpolation
and an eventual path to posterior uncertainty is preferred over a neural
closure. Use Koopman-EDMD, subspace, or other lifted/input-output surrogates
when predictive accuracy over the tested envelope is more important than
recovering aerodynamic parameters. Use local model stitching or LPV-style
scheduling when separate parts of the flight envelope are well represented by
simple models but one global model averages incompatible trim, stall-onset, and
recovery behavior. Use
hidden-input models when a proprietary autopilot reshapes commands. Use
frequency-domain tools when the experiment has been designed for spectral
identification. The value of the benchmark is that each method is exposed to
the same train/validation splits across these regimes, so these recommendations
can be revised quantitatively as additional methods and datasets are added.

The updated aggressive results also clarify the role of simple linear models.
Linear-SS is not invalid; it is an effective local predictor for small motions
and can remain competitive when mocap-derived state reconstruction dominates
the error budget. It should not, however, be interpreted as recovering the
nonlinear aerodynamics. When direct states expose the stall and recovery
regimes, the nonlinear closure and residual-learning methods give substantially
lower validation error. This is the intended benchmark behavior: near-trim tests
verify that simple models are not unfairly penalized, while aggressive
initial-condition and maneuver variation reveal when global linearity is no
longer adequate.

\section{Repository Roadmap}
The immediate roadmap is repository-driven rather than submission-driven. The first priority is to keep the 3-DOF benchmark reproducible: generated datasets, figures, tables, website artifacts, and this reference document should be regenerated from the same command-line entry points. The second priority is to mature the 6-DOF nonlinear benchmark methods beyond the initial attached-flow, linear, residual, and mocap-output baselines. A useful future extension is a separate F-16 model family based on the Stevens--Lewis table-lookup aerodynamics, which would test high-performance aircraft identification without replacing the present small-RC, mocap-motivated benchmark. The third priority is to stabilize the method-plugin interface so external contributors can add methods without editing the core comparison harness.

Future methods should be added only when the repository contains executable code, benchmark metadata, and validation results for the new row. Candidate additions include input-output nonlinear predictors, richer local/LPV model stitching, Bayesian uncertainty-aware variants of the existing grey-box methods, and high-fidelity aerodynamic surrogate closures. Until those implementations exist in the repository, this reference limits detailed method descriptions to the rows already represented by code and generated results.

\appendix
\section{Benchmark Result Tables}
\label{app:benchmark_tables}
The main text uses figures to summarize trends. The following tables retain the
complete sorted numerical results used to generate those figures. Direct-state
tables use simulated noisy $V,\alpha,\gamma,Q$ measurements. Mocap-derived
tables use states reconstructed from 100 Hz position and pitch-attitude
measurements. Lower validation score is better.

\subsection{6-DOF Baseline Results}
\IfFileExists{generated/aircraft6dof_method_comparison_table.tex}{%
\input{generated/aircraft6dof_method_comparison_table}
}{}

\IfFileExists{generated/shared_method_comparison_direct_table.tex}{%
\input{generated/shared_method_comparison_direct_table}
}{}
\IfFileExists{generated/shared_method_comparison_mocap_table.tex}{%
\input{generated/shared_method_comparison_mocap_table}
}{}

\subsection{SAFE-Off and SAFE-On Open-Loop Cases}
\IfFileExists{generated/open_loop_shared_method_comparison_direct_table.tex}{%
\input{generated/open_loop_shared_method_comparison_direct_table}
}{}
\IfFileExists{generated/open_loop_shared_method_comparison_mocap_table.tex}{%
\input{generated/open_loop_shared_method_comparison_mocap_table}
}{}
\IfFileExists{generated/open_loop_safe_shared_method_comparison_direct_table.tex}{%
\input{generated/open_loop_safe_shared_method_comparison_direct_table}
}{}
\IfFileExists{generated/open_loop_safe_shared_method_comparison_mocap_table.tex}{%
\input{generated/open_loop_safe_shared_method_comparison_mocap_table}
}{}

\subsection{SAFE Recovery-Probe Case}
\IfFileExists{generated/safe_loop_shared_method_comparison_direct_table.tex}{%
\input{generated/safe_loop_shared_method_comparison_direct_table}
}{}
\IfFileExists{generated/safe_loop_shared_method_comparison_mocap_table.tex}{%
\input{generated/safe_loop_shared_method_comparison_mocap_table}
}{}

\subsection{Frequency-Excitation Cases}
\IfFileExists{generated/sine_sweep_shared_method_comparison_direct_table.tex}{%
\input{generated/sine_sweep_shared_method_comparison_direct_table}
}{}
\IfFileExists{generated/sine_sweep_shared_method_comparison_mocap_table.tex}{%
\input{generated/sine_sweep_shared_method_comparison_mocap_table}
}{}
\IfFileExists{generated/sine_sweep_safe_shared_method_comparison_direct_table.tex}{%
\input{generated/sine_sweep_safe_shared_method_comparison_direct_table}
}{}
\IfFileExists{generated/sine_sweep_safe_shared_method_comparison_mocap_table.tex}{%
\input{generated/sine_sweep_safe_shared_method_comparison_mocap_table}
}{}

\subsection{Aggressive Nonlinear Maneuver Cases}
\IfFileExists{generated/aggressive_shared_method_comparison_direct_table.tex}{%
\input{generated/aggressive_shared_method_comparison_direct_table}
}{}
\IfFileExists{generated/aggressive_shared_method_comparison_mocap_table.tex}{%
\input{generated/aggressive_shared_method_comparison_mocap_table}
}{}
\IfFileExists{generated/aggressive_safe_shared_method_comparison_direct_table.tex}{%
\input{generated/aggressive_safe_shared_method_comparison_direct_table}
}{}
\IfFileExists{generated/aggressive_safe_shared_method_comparison_mocap_table.tex}{%
\input{generated/aggressive_safe_shared_method_comparison_mocap_table}
}{}

\section{Observation-Rate and UQ Tables}
\IfFileExists{generated/observation_rate_table.tex}{%
\input{generated/observation_rate_table}
}{}
\IfFileExists{generated/shared_uq_diagnostics_table.tex}{%
\input{generated/shared_uq_diagnostics_table}
}{}

\bibliographystyle{plainnat}
\bibliography{ref}


\end{document}
